\documentclass[10pt,a4paper,landscape]{report}

\usepackage[utf8]{inputenc}
\usepackage{amsmath}
\usepackage{amsfonts}
\usepackage{amssymb}
\usepackage{graphicx}
\usepackage{xcolor}
\definecolor{gris}{rgb}{0.75,0.75,0.75}
\usepackage{colortbl}
\usepackage{enumitem}
\setitemize{noitemsep,topsep=0pt,parsep=0pt,partopsep=0pt,leftmargin=5pt}
\usepackage{version}
\usepackage{multirow}
\usepackage[top=1cm, bottom=2cm, left=2cm, right=2cm]{geometry}

\definecolor{true}{rgb}{0,0.40,0}
\definecolor{ukwn}{rgb}{0,0,0.50}
\definecolor{false}{rgb}{0.60,0,0}

\usepackage{fancyhdr}
\pagestyle{fancy}

\fancyhead{}
\fancyfoot{} 
\lhead{\includegraphics{../Logo/logoKNKMini.jpg} \hspace{0.1cm} Kould Not Konect  \hspace{0.4cm} \vline}
\chead{Cahier des tests de recette}
\rhead{Kould Not Share}
\rfoot{\thepage}

\author{Nicolas REYNAUD, Kevin BASCOL, Kévin Laoussing}
\title{Cahier des tests de recette}
\date{5 Novembre 2014}

\makeatletter
\renewcommand{\thesection}{\@arabic\c@section}
\makeatother

\makeatletter
\def\chline#1{%
  \noalign{\ifnum0=`}\fi\begingroup\color{#1}\hrule \@height \arrayrulewidth\endgroup \futurelet
   \reserved@a\@xhline}
\def\ccline#1#2{\@ccline#1\@nil#2\@nil}
\def\@ccline#1-#2\@nil#3\@nil{%
  \omit
  \@multicnt#1%
  \advance\@multispan\m@ne
  \ifnum\@multicnt=\@ne\@firstofone{&\omit}\fi
  \@multicnt#2%
  \advance\@multicnt-#1%
  \advance\@multispan\@ne
  \begingroup\color{#3}\leaders\hrule\@height\arrayrulewidth\hfill\endgroup
  \cr
  \noalign{\vskip-\arrayrulewidth}}
\makeatother

\begin{document}
\makeatletter
	\begin{titlepage}
	
	\begin{figure}
		\begin{minipage}[c]{.46\linewidth}
		\end{minipage} \hfill
		\begin{minipage}[c]{.20\linewidth}
			\begin{center}
				\includegraphics{../Logo/logoKNK.jpg}\\
				{\large Kould Not Konect}
			\end{center}
		\end{minipage}
	\vspace{1cm}
	\end{figure}
	
	\centering
		{
		\hrule height 2pt
		\vspace{0.7cm}
		\Huge \textbf{\@title}}\\
		\vspace{0.7cm}
		\hrule height 2pt
		\vspace{1.5cm}
		{\LARGE  Projet \textbf{Kould Not Share} v1.0}
		
		\vfill
		
		\begin{tabular}{|c|c|c|}
			\hline
			Version & Date & Description\\
			\hline
			v.1 & 05/11/14 & Tests de la première version des exigences\\
			\hline
			 & &\\
			\hline
			 & &\\
			\hline
		\end{tabular}\\
		\vspace{1cm}
		\@author\\
		\end{titlepage}
\makeatother
\setcounter{secnumdepth}{5}
\setcounter{tocdepth}{5}
\renewcommand{\contentsname}{Sommaire}
\begingroup\makeatletter
\def\@makeschapterhead#1{%
  {\parindent \z@ \raggedright
    \normalfont
    \interlinepenalty\@M
    \Huge \bfseries  #1\par\nobreak
    \vskip 20pt% <---- à réduire pour avoir plus de place
  }}\makeatother
\tableofcontents
\endgroup
\thispagestyle{empty}
\setcounter{page}{0}
\newpage

\newgeometry{top=2cm, bottom=2cm, left=0.5cm, right=0.5cm}

\section{Tests Client}
\subsection{Authentification}
\begin{center}
	\bgroup
	\def\arraystretch{1.5}
	\begin{tabular}{|p{2.5cm}|p{2cm}|p{8cm}|p{8cm}|p{5cm}|}
		\hline
		\rowcolor{gris}Référence & Exigence & Test & Résultat attendu & Résultat/Commentaires\\
		\hline
		AuthC.1 & 3.1.1.1.1.1 3.1.1.1.1.2 3.1.1.1.1.3 3.1.1.1.1.13 & Ouvrir le logiciel. & Le formulaire de connexion apparait en pop-up au centre de la fenêtre, le reste de l'application est grisé (hormis la barre d'action). &\vspace*{-0.2cm} \begin{itemize}[label=$\ast$] \item \textcolor{true}{Le pop-up est apparu au centre.} \item \textcolor{false}{Le reste de l'application n'est pas grisé.}\end{itemize}\\
		\cline{1-2}\ccline{3-3}{gris}\cline{4-5}
		AuthC.1.1 & 3.1.1.1.1.4 3.1.1.1.1.5 & Cliquer sur le bouton de déconnexion. & Le bouton doit être grisé et ne doit rien déclencher. &\vspace*{-0.2cm} \begin{itemize}[label=$\ast$] \item \textcolor{true}{Le bouton est désactivé.} \item \textcolor{false}{Le bouton n'est pas grisé.}\end{itemize}\\
		\cline{1-2}\ccline{3-3}{gris}\cline{4-5}
		AuthC.1.2 & 3.1.1.1.1.6 3.1.1.1.1.7 3.1.1.1.1.8 &  & sur le pop-up, le titre "Authentification" apparait centré et est au dessus d'une barre horizontale. &\vspace*{-0.2cm} \begin{itemize}[label=$\ast$] \item \textcolor{true}{Validé.}\end{itemize}\\
		\cline{1-2}\cline{4-5}
		AuthC.1.3 & 3.1.1.1.1.9 3.1.1.1.1.10 3.1.1.1.2.3 3.1.1.1.2.4 & & Les mots "Serveur", "Nom d'utilisateur" et "Mot de passe" sont affichés, chacun suivit d'un champ de saisie de texte. &\vspace*{-0.2cm} \begin{itemize}[label=$\ast$] \item \textcolor{true}{Validé.}\end{itemize}\\
		\cline{1-2}\cline{4-5}
		AuthC.1.4 & 3.1.1.1.1.11 & & Un bouton "Se connecter" est affiché centré en bas du pop-up &\vspace*{-0.2cm}  \begin{itemize}[label=$\ast$] \item \textcolor{true}{Validé.}\end{itemize}\\
		\cline{1-2}\ccline{3-3}{gris}\cline{4-5}
		AuthC.1.4.1 & 3.1.1.1.1.12 3.1.1.1.2.8 3.1.1.1.2.9 3.1.1.1.2.10 & Remplir incorrectement un ou plusieurs champs de saisie puis cliquer sur le bouton "Se connecter" (à effectuer avec plusieurs combinaison de champs incorrects). & Le(s) champ(s) incorrect(s) possède(nt) une bordure rouge. &\vspace*{-0.2cm} \begin{itemize}[label=$\ast$] \item \textcolor{false}{Non impémenté}\end{itemize}\\
		\cline{1-2}\ccline{3-3}{gris}\cline{4-5}
		AuthC.1.4.1.1 & 3.1.1.1.2.11 & Rentrer une adresse serveur autre qu'une adresse IP v4. & La bordure du champ "Serveur" devient rouge et une boite de dialogue avec le texte "Erreur de serveur". &\vspace*{-0.2cm} \begin{itemize}[label=$\ast$] \item \textcolor{false}{Non implémenté}\end{itemize}\\
		\hline
	\end{tabular}
	\egroup
\end{center}
\newpage
\begin{center}
	\bgroup
	\def\arraystretch{1.5}
	\begin{tabular}{|p{2.5cm}|p{2cm}|p{8cm}|p{8cm}|p{5cm}|}
		\hline
		\rowcolor{gris}Référence & Exigence & Test & Résultat attendu & Résultat/Commentaires\\
		\hline
		AuthC.1.4.1.2 & 3.1.1.1.2.12 3.1.1.1.2.15.2 & Rentrer un identifiant de plus de 30 caractères. & La bordure du champ "Nom d'utilisateur" devient rouge et une boite de dialogue avec le texte "L'utilisateur entré est invalide". &\vspace*{-0.2cm} \begin{itemize}[label=$\ast$] \item  \textcolor{true}{L'erreur "L'utilisateur entré est invalide" s'affiche en cas de pseudo trop long} \item \textcolor{false}{Le champ ne devient pas rouge.} \item \textcolor{false}{L'erreur est affichée directement dans le pop-up}\end{itemize}\\
		\cline{1-2}\ccline{3-3}{gris}\cline{4-5}
		AuthC.1.4.1.3 & 3.1.1.1.2.13 3.1.1.1.2.15.2 & Rentrer un identifiant de moins de 3 caractères. & La bordure du champ "Nom d'utilisateur" devient rouge et une boite de dialogue avec le texte "L'utilisateur entré est invalide". &\vspace*{-0.2cm} \begin{itemize}[label=$\ast$] \item \textcolor{false}{Non impémenté}\end{itemize}\\
		\cline{1-2}\ccline{3-3}{gris}\cline{4-5}
		AuthC.1.4.1.4 & 3.1.1.1.2.14 3.1.1.1.2.15.1 & Rentrer un mot de passe de moins de 6 caractères. & La bordure du champ "Mot de passe" devient rouge et une boite de dialogue avec le texte "Erreur de mot de passe". &\vspace*{-0.2cm} \begin{itemize}[label=$\ast$] \item  \textcolor{true}{L'erreur s'affiche en cas de mot de passe trop court} \item \textcolor{false}{Le champ ne devient pas rouge.} \item \textcolor{false}{L'erreur est affichée directement dans le pop-up}\end{itemize}\\
		
		\cline{1-2}\ccline{3-3}{gris}\cline{4-5}
		AuthC.1.4.1.5 & 3.1.1.1.2.15 3.1.1.1.2.15.1 & Rentrer un mot de passe de plus de 50 caractères. & La bordure du champ "Mot de passe" devient rouge et une boite de dialogue avec le texte "Erreur de mot de passe". &\vspace*{-0.2cm} \begin{itemize}[label=$\ast$] \item  \textcolor{true}{L'erreur s'affiche en cas de mot de passe trop long} \item \textcolor{false}{Le champ ne devient pas rouge.} \item \textcolor{false}{L'erreur est affichée directement dans le pop-up}\end{itemize}\\
		\cline{1-2}\ccline{3-3}{gris}\cline{4-5}
		AuthC.1.4.2 & 3.1.1.1.2.1 3.1.1.1.2.5 3.1.1.1.2.6 & Remplir correctement les champs, puis cliquer sur "Se connecter". & Le programme se connecte au serveur avec l'identifiant donné. &\vspace*{-0.2cm} \begin{itemize}[label=$\ast$] \item \textcolor{true}{Validé.}\end{itemize}\\
		\cline{1-2}\cline{4-5}
		AuthC.1.4.2.1 & 3.1.1.1.2.7 & & Le message "Bonjour [pseudo]" apparait dans la barre de menu. &\vspace*{-0.2cm} \begin{itemize}[label=$\ast$] \item \textcolor{false}{Non impémenté}\end{itemize}\\
		\cline{1-2}\cline{4-5}
		AuthC.1.4.2.2 & 3.1.1.1.2.16 & & Une case à cocher est affiché à côté du message "Bonjour [pseudo]". &\vspace*{-0.2cm} \begin{itemize}[label=$\ast$] \item \textcolor{false}{Non impémenté}\end{itemize}\\
		\cline{1-2}\ccline{3-3}{gris}\cline{4-5}
		AuthC.1.4.2.2.1 & 3.1.1.1.2.16 & Cocher la case, fermer le programme et le rouvrir.& Le programme se connecte automatiquement au serveur précédent. &\vspace*{-0.2cm} \begin{itemize}[label=$\ast$] \item \textcolor{false}{Non impémenté}\end{itemize}
\\
		\hline
		AuthC.1.5 & 3.1.1.1.2.2 & & Le programme ne propose aucun moyen de créer ou modifier des identifiants et mots de passe. &\vspace*{-0.2cm} \begin{itemize}[label=$\ast$] \item \textcolor{true}{Validé.}\end{itemize}\\
		\hline
	\end{tabular}
	\egroup
\end{center}

\subsection{Gestionnaire de marques-page}
\begin{center}
	\bgroup
	\def\arraystretch{1.5}
	\begin{tabular}{|p{2.5cm}|p{2cm}|p{8cm}|p{8cm}|p{5cm}|}
		\hline
		\rowcolor{gris}Référence & Exigence & Test & Résultat attendu & Résultat/Commentaires\\
		\hline
		GdmpC.1 & 3.1.1.2.1 3.1.1.2.2 & Ouvrir le logiciel. & Un menu déroulant "Bookmarks" est présent dans la barre d'action. &\vspace*{-0.2cm} \begin{itemize}[label=$\ast$] \item \textcolor{false}{Non impémenté}\end{itemize}\\
		\cline{1-2}\ccline{3-3}{gris}\cline{4-5}
		GdmpC.1.1 & 3.1.1.2.3 3.1.1.2.7& Dérouler le menu "Bookmarks". & Les boutons "Afficher les marque-pages" et "Mettre un marque-page" sont présents dans le menu. &\vspace*{-0.2cm} \begin{itemize}[label=$\ast$] \item \textcolor{false}{Non impémenté}\end{itemize} \\
		\cline{1-2}\ccline{3-3}{gris}\cline{4-5}
		GdmpC.1.1.1 & 3.1.1.2.4 & Cliquer sur "Afficher les marque-pages". & un pop-up s'ouvre avec la liste des liens vers les serveurs affichés avec les noms des marque-pages &\vspace*{-0.2cm} \begin{itemize}[label=$\ast$] \item \textcolor{false}{Non impémenté}\end{itemize}
\\
		\cline{1-2}\ccline{3-3}{gris}\cline{4-5}
		GdmpC.1.1.2 & 3.1.1.2.8 & Cliquer sur "Mettre un marque-page". & Un champs de saisie "nom" et un case à cocher "enregistrer mot de passe" apparaissent dans une boite de dialogue. &\vspace*{-0.2cm} \begin{itemize}[label=$\ast$] \item \textcolor{false}{Non impémenté}\end{itemize}
\\
		\cline{1-2}\ccline{3-3}{gris}\cline{4-5}
		GdmpC.1.1.2.1 & 3.1.1.2.9 & Entrer un nom. & Le marque page est ajouté à la liste. &\vspace*{-0.2cm} \begin{itemize}[label=$\ast$] \item \textcolor{false}{Non impémenté}\end{itemize}
\\
		\cline{1-2}\ccline{3-3}{gris}\cline{4-5}
		GdmpC.1.1.2.2 & 3.1.1.2.10 & Entrer un nom et cocher la case. & Le marque-page est ajouté à la liste et permet une connexion directe. &\vspace*{-0.2cm} \begin{itemize}[label=$\ast$] \item \textcolor{false}{Non impémenté}\end{itemize}
\\
		\hline
		GdmpC.2 & 3.1.1.2.5 & Ouvrir le fichier des marque-pages. & Les marque-pages sont composés d'un nom, de l'adresse du serveur et de l'identifiant de l'utilisateur. &\vspace*{-0.2cm} \begin{itemize}[label=$\ast$] \item \textcolor{false}{Non impémenté}\end{itemize}
\\
		\cline{1-2}\cline{4-5}
		GdmpC.2.1 & 3.1.1.2.6 & & Les marque-pages où le mot de passe est enregistré sont composé aussi du mot de passe crypté. &\vspace*{-0.2cm} \begin{itemize}[label=$\ast$] \item \textcolor{false}{Non impémenté}\end{itemize}
\\
		\hline
	\end{tabular}
	\egroup
\end{center}

\subsection{Limitations}
\begin{center}
	\bgroup
	\def\arraystretch{1.5}
	\begin{tabular}{|p{2.5cm}|p{2cm}|p{8cm}|p{8cm}|p{5cm}|}
		\hline
		\rowcolor{gris}Référence & Exigence & Test & Résultat attendu & Résultat/Commentaires\\
		\hline
		LimiC.1 & 3.1.1.3.1 & Ouvrir le logiciel. & Le menu déroulant "Configuration est présent dans la barre d'action. &\vspace*{-0.2cm} \begin{itemize}[label=$\ast$] \item \textcolor{false}{Non impémenté}\end{itemize}\\
		\cline{1-2}\ccline{3-3}{gris}\cline{4-5}
		LimiC.1.1 & 3.1.1.3.3 3.1.1.3.7 3.1.1.3.11 3.1.1.3.14 & Cliquer sur le menu déroulant "Configuration". & Les boutons "Limiter vitesse téléchargement", "Limiter vitesse téléversement", "Plages horaires téléchargement" et "Plages horaires téléversement" sont présent dans le menu déroulant. &\vspace*{-0.2cm} \begin{itemize}[label=$\ast$] \item \textcolor{false}{Non impémenté}\end{itemize}
\\
		\cline{1-2}\ccline{3-3}{gris}\cline{4-5}
		LimiC.1.1.1 & 3.1.1.3.4 3.1.1.3.5 & Cliquer sur "Limiter vitesse téléchargement". & Une boite de dialogue avec un curseur apparait. &\vspace*{-0.2cm} \begin{itemize}[label=$\ast$] \item \textcolor{false}{Non impémenté}\end{itemize}
\\
		\cline{1-2}\ccline{3-3}{gris}\cline{4-5}
		LimiC.1.1.1.2 & 3.1.1.3.2 & Choisir une limite et valider la boite de dialogue. & La vitesse de téléchargement ne dépasse pas la limite fixée. &\vspace*{-0.2cm} \begin{itemize}[label=$\ast$] \item \textcolor{false}{Non impémenté}\end{itemize}
\\
		\cline{1-2}\ccline{3-3}{gris}\cline{4-5}
		LimiC.1.1.2 & 3.1.1.3.8 3.1.1.3.9 & Cliquer sur "Limiter vitesse téléversement". & Une boite de dialogue avec un curseur apparait. &\vspace*{-0.2cm} \begin{itemize}[label=$\ast$] \item \textcolor{false}{Non impémenté}\end{itemize}
\\
		\cline{1-2}\ccline{3-3}{gris}\cline{4-5}
		LimiC.1.1.2.1 & 3.1.1.3.6 & Choisir une limite et valider la boite de dialogue. & La vitesse de téléversement ne dépasse pas la limite fixée. &\vspace*{-0.2cm} \begin{itemize}[label=$\ast$] \item \textcolor{false}{Non impémenté}\end{itemize}
\\
		\cline{1-2}\ccline{3-3}{gris}\cline{4-5}
		LimiC.1.1.3 & 3.1.1.3.12 & Cliquer sur "Plages horaires téléchargement". & Une boite de dialogue apparait avec la liste des heures du jours couplées à une case à cocher. &\vspace*{-0.2cm} \begin{itemize}[label=$\ast$] \item \textcolor{false}{Non impémenté}\end{itemize}
\\
		\cline{1-2}\ccline{3-3}{gris}\cline{4-5}
		LimiC.1.1.3.1 & 3.1.1.3.10 & Sélectionner des horaires et valider la boite de dialogue. & Aux heures sélectionnées le programme bloque le téléchargement. &\vspace*{-0.2cm} \begin{itemize}[label=$\ast$] \item \textcolor{false}{Non impémenté}\end{itemize}
\\
		\cline{1-2}\ccline{3-3}{gris}\cline{4-5}
		LimiC.1.1.4 & 3.1.1.3.15 & Cliquer sur "Plages horaires téléversement". & Une boite de dialogue apparait avec la liste des heures du jours couplées à une case à cocher. &\vspace*{-0.2cm} \begin{itemize}[label=$\ast$] \item \textcolor{false}{Non impémenté}\end{itemize}
\\
		\cline{1-2}\ccline{3-3}{gris}\cline{4-5}
		LimiC.1.1.4.1 & 3.1.1.3.13 & Sélectionner des horaires et valider la boite de dialogue. & Aux heures sélectionnées le programme bloque le téléversement. &\vspace*{-0.2cm} \begin{itemize}[label=$\ast$] \item \textcolor{false}{Non impémenté}\end{itemize}
\\
		\hline
	\end{tabular}
	\egroup
\end{center}

\subsection{Utilisation des dossiers}
\begin{center}
	\bgroup
	\def\arraystretch{1.5}
	\begin{tabular}{|p{2.5cm}|p{2cm}|p{8cm}|p{8cm}|p{5cm}|}
		\hline
		\rowcolor{gris}Référence & Exigence & Test & Résultat attendu & Résultat/Commentaires\\
		\hline
		UtdoC.1 & 3.1.1.4.1 & Ouvrir le logiciel. & L'arborescence des dossiers dont l'utilisateur a accès sur le serveur est affichée sur la droite de la fenêtre. &\vspace*{-0.2cm} \begin{itemize}[label=$\ast$] \item \textcolor{true}{Validé.}\end{itemize} \\
		\cline{1-2}\ccline{3-3}{gris}\cline{4-5}
		UtdoC.1.1 & 3.1.1.4.2 3.1.1.4.3 & Cliquer sur les dossiers dans l'arborescence. & Les sous-dossiers et fichiers des dossiers sur lesquels on a cliqué sont affichés dans l'arborescence. &\vspace*{-0.2cm} \begin{itemize}[label=$\ast$] \item \textcolor{true}{Validé.}\end{itemize} \\
		\cline{1-2}\ccline{3-3}{gris}\cline{4-5}
		UtdoC.1.2 & 3.1.1.4.5 3.1.1.4.6 3.1.1.4.10 3.1.1.4.13 & Faire un clic droit sur un dossier. & Un menu de gestion de dossier apparait, les boutons "Nouveau sous-dossier", "Supprimer le dossier", et "Nouveau fichier vide" sont disponible à l'intérieur. & \vspace*{-0.2cm} \begin{itemize}[label=$\ast$] \item \textcolor{false}{Non impémenté}\end{itemize} \\
		\cline{1-2}\ccline{3-3}{gris}\cline{4-5}
		UtdoC.1.2.1 & 3.1.1.4.7 3.1.1.4.8 & Cliquer sur "Nouveau sous-dossier". & Une boite de dialogue apparait avec un champ de saisie "nom". &\vspace*{-0.2cm} \begin{itemize}[label=$\ast$] \item \textcolor{false}{Non impémenté}\end{itemize}
\\
		\cline{1-2}\ccline{3-3}{gris}\cline{4-5}
		UtdoC.1.2.1.1 & 3.1.1.4.4 & Entrer un nom et valider la boite de dialogue. & Un nouveau sous-dossier avec le nom donné dans la boite de dialogue apparait dans l'arborescence. &\vspace*{-0.2cm} \begin{itemize}[label=$\ast$] \item \textcolor{false}{Non impémenté}\end{itemize}
\\
		\cline{1-2}\ccline{3-3}{gris}\cline{4-5}
		UtdoC.1.2.2 & 3.1.1.4.11 3.1.1.4.11.1 & Cliquer sur "Supprimer le dossier". & Une boite de dialogue proposant un bouton "Valider" et un "Annuler" apparait. &\vspace*{-0.2cm} \begin{itemize}[label=$\ast$] \item \textcolor{false}{Non impémenté}\end{itemize}
\\
		\cline{1-2}\ccline{3-3}{gris}\cline{4-5}
		UtdoC.1.2.2.1 & 3.1.1.4.9 3.1.1.4.11.2 & Cliquer sur "Valider". & La boite de dialogue disparait et le dossier est supprimé. &\vspace*{-0.2cm} \begin{itemize}[label=$\ast$] \item \textcolor{false}{Non impémenté}\end{itemize}
\\
		\cline{1-2}\ccline{3-3}{gris}\cline{4-5}
		UtdoC.1.2.2.2 & 3.1.1.4.11.3 & Cliquer sur "Annuler". & La boite de dialogue disparait sans rien déclencher. &\vspace*{-0.2cm} \begin{itemize}[label=$\ast$] \item \textcolor{false}{Non impémenté}\end{itemize}
\\
		\cline{1-2}\ccline{3-3}{gris}\cline{4-5}
		UtdoC.1.2.3 & 3.1.1.4.12 3.1.1.4.14 & Cliquer sur "Nouveau fichier vide". & Un nouveau fichier "nouveau\_fichier" apparait dans le dossier. &\vspace*{-0.2cm} \begin{itemize}[label=$\ast$] \item \textcolor{false}{Non impémenté}\end{itemize}
\\
		\cline{1-2}\ccline{3-3}{gris}\cline{4-5}
		Utdoc.1.3 & 3.1.1.4.16 3.1.1.4.17 & Faire un clic droit sur un des fichiers dans l'arborescence. & Un menu de gestion de fichier apparait avec un bouton "Supprimer le fichier" à l'intérieur. &\vspace*{-0.2cm} \begin{itemize}[label=$\ast$] \item \textcolor{false}{Non impémenté}\end{itemize} \\
		\cline{1-2}\ccline{3-3}{gris}\cline{4-5}
		UtdoC.1.3.2 & 3.1.1.4.18 3.1.1.4.18.1 & Cliquer sur "Supprimer le fichier". & Une boite de dialogue proposant un bouton "Valider" et un "Annuler" apparait. &\vspace*{-0.2cm} \begin{itemize}[label=$\ast$] \item \textcolor{false}{Non impémenté}\end{itemize}
\\
		\cline{1-2}\ccline{3-3}{gris}\cline{4-5}
		UtdoC.1.3.2.1 & 3.1.1.4.15 3.1.1.4.18.2 & Cliquer sur "Valider". & La boite de dialogue disparait et le dossier est supprimé. &\vspace*{-0.2cm} \begin{itemize}[label=$\ast$] \item \textcolor{false}{Non impémenté}\end{itemize}
\\
		\cline{1-2}\ccline{3-3}{gris}\cline{4-5}
		UtdoC.1.3.2.2 & 3.1.1.4.18.3 & Cliquer sur "Annuler". & La boite de dialogue disparait sans rien déclencher. &\vspace*{-0.2cm} \begin{itemize}[label=$\ast$] \item \textcolor{false}{Non impémenté}\end{itemize}
\\
		\hline
	\end{tabular}
	\egroup
\end{center}

\subsection{Communication réseau}

\subsubsection{Procédure d'initialisation de la communication}

\begin{center}
	\bgroup
	\def\arraystretch{1.5}
	\begin{tabular}{|p{2.5cm}|p{2cm}|p{8cm}|p{8cm}|p{5cm}|}
		\hline
		\rowcolor{gris}Référence & Exigence & Test & Résultat attendu & Résultat/Commentaires\\
		\hline
		PicoC.1.1 & 3.1.1.5.1.1 & Après la connexion du client vers le serveur, lancer la commande "netstat -a -p" sur le terminal du client & Le terminal doit afficher une connexion avec une adresse IP correspondant au serveur, dont le port de connexion doit être 21, l'état doit être à "ESTABLISHED" et le nom ou le PID doit être celui du logiciel.&\vspace*{-0.2cm} \begin{itemize}[label=$\ast$] \item \textcolor{true}{Validé.}\end{itemize} \\
		\hline
		PicoC.1.2 & 3.1.1.5.1.2 3.1.1.5.1.4 & Après le lancement d'un transfert de fichier, lancer la commande "netstat -a -p" sur le terminal du client & Le terminal doit afficher une connexion avec une adresse IP correspondant au serveur, dont le port de connexion doit être différent de 21 (soit supérieur à 1024, soit égale à 20), l'état doit être à "ESTABLISHED" et le nom ou le PID doit être celui du logiciel. &\vspace*{-0.2cm} \begin{itemize}[label=$\ast$] \item \textcolor{true}{Validé.}\end{itemize} \\
		\hline
		PicoC.1.3 & 3.1.1.5.1.3 & Après le lancement d'un transfert de fichier, lancer la commande "netstat -a" sur le terminal du client & Le terminal ne doit plus afficher une connexion avec une adresse IP correspondant au serveur \textbf{et} dont le port de connexion est 21. &\vspace*{-0.2cm} \begin{itemize}[label=$\ast$] \item \textcolor{false}{Non implémenté.}\end{itemize}\\
		\hline
	\end{tabular}
	\egroup
\end{center}

\subsubsection{Procédure en cas d'échec d'initialisation de la communication}

\begin{center}
	\bgroup
	\def\arraystretch{1.5}
	\begin{tabular}{|p{2.5cm}|p{2cm}|p{8cm}|p{8cm}|p{5cm}|}
		\hline
		\rowcolor{gris}Référence & Exigence & Test & Résultat attendu & Résultat/Commentaires\\
		\hline
		PecoC.1.1.a & 3.1.1.5.2.1.1.a & Débrancher la machine contenant le logiciel client du réseau, et faire une demande de connexion avec un nom de serveur FTP valide.& Une fenêtre de dialogue doit s'ouvrir avec le message suivant :   "Impossible d'établir une connexion avec le serveur $< nom\_du\_serveur>$:\linebreak
-Vérifiez l'état de votre connexion Internet.\linebreak
-Vérifiez que le nom du serveur n'est pas erroné.".&\vspace*{-0.2cm} \begin{itemize}[label=$\ast$] \item \textcolor{false}{Non implémenté.}\end{itemize} \\
		\cline{1-3}\cline{5-5}
		PecoC.1.1.b & 3.1.1.5.2.1.1.b & Faire une demande de connexion avec un nom de serveur FTP non-valide.& .&\vspace*{-0.2cm} \begin{itemize}[label=$\ast$] \item \textcolor{false}{Non implémenté.}\end{itemize}\\
		\hline
		PecoC.1.2 & 3.1.1.5.2.1.2 & Faire une demande de connexion avec un nom de serveur FTP valide dont le nombre de connexion autorisé par le serveur est atteint.& Une fenêtre de dialogue doit s'ouvrir avec le message suivant :   "Impossible d'établir une connexion avec le serveur $< nom\_du\_serveur>$: le serveur est surchargé.".&\vspace*{-0.2cm} \begin{itemize}[label=$\ast$] \item \textcolor{false}{Non implémenté.}\end{itemize}\\
		\hline
		PecoC.1.3  & 3.1.1.5.2.1.3 & Faire une demande de connexion avec un nom de serveur FTP valide. Le serveur doit avoir auparavant banni le client du test.& Une fenêtre de dialogue doit s'ouvrir avec le message suivant :   "Impossible d'établir une communication avec le serveur $< nom\_du\_serveur>$:
$< nom\_du\_serveur>$ vous a banni !".&\vspace*{-0.2cm} \begin{itemize}[label=$\ast$] \item \textcolor{false}{Non implémenté.}\end{itemize}\\
		\hline
		PecoC.2.2 & 3.1.1.5.2.2 & Faire une demande de connexion avec un nom de serveur FTP valide. Les regles du pare-feu doit être configurer pour que toutes les connexions sortant vers des ports supérieurs à 1024 ne soient pas autorisées.& Une fenêtre de dialogue doit s'ouvrir avec le message suivant :   "Impossible d'établir une connexion avec le serveur $< nom\_du\_serveur>$ sur le port $< numero\_du\_port\_echange>$ : Vérifiez vos règle de pare-feu.".&\vspace*{-0.2cm} \begin{itemize}[label=$\ast$] \item \textcolor{false}{Non implémenté.}\end{itemize}\\
		\hline
		PecoC.2.3 & 3.1.1.5.2.3 & Vérifier pour les tests de la section 3.1.1.5 que la fenêtre possede bien un bouton "Fermer" fonctionnel :  cliquer sur ce bouton doit fermer la fenêtre.& Cliquer sur le bouton "Fermer" provoque le fermeture de la fenêtre de dialogue .&\vspace*{-0.2cm} \begin{itemize}[label=$\ast$] \item \textcolor{false}{Non implémenté.}\end{itemize}\\
		\hline
		PecoC.2.4 & 3.1.1.5.2.4 & Vérifier pour les tests de la section PecoC.2.* que la fenêtre possede bien un bouton "Ré-essayer" fonctionnel :  cliquer sur ce bouton doit ré-itérer l'opération échoué.& Cliquer sur le bouton "Ré-essayer" provoque la ré-itération de l'opération échoué.&\vspace*{-0.2cm} \begin{itemize}[label=$\ast$] \item \textcolor{false}{Non implémenté.}\end{itemize}\\
		\hline
	\end{tabular}
	\egroup
\end{center}

\subsubsection{Procédure en cas de coupure de la communication}

\begin{center}
	\bgroup
	\def\arraystretch{1.5}
	\begin{tabular}{|p{2.5cm}|p{2cm}|p{8cm}|p{8cm}|p{5cm}|}
		\hline
		\rowcolor{gris}Référence & Exigence & Test & Résultat attendu & Résultat/Commentaires\\
		\hline
		PecoC.1 & 3.1.1.5.3.1 & Pendant un téléchargement, couper la ligne internet de la machine cliente & Une fenêtre de dialogue doit s'ouvrir avec le message suivant :   "Connexion interrompue avec le serveur $< nom\_du\_serveur>$.".&\vspace*{-0.2cm} \begin{itemize}[label=$\ast$] \item \textcolor{false}{Non implémenté.}\end{itemize} \\
		\hline
		PccoC.2 & 3.1.1.5.3.2 & Idem que le test PecoC.2.3. & Idem que le test PecoC.2.3. & \vspace*{-0.2cm} \begin{itemize}[label=$\ast$] \item \textcolor{false}{Non implémenté.}\end{itemize}\\
		\hline
	\end{tabular}
	\egroup
\end{center}

\subsubsection{Procédure en cas de fermeture de la communication}

\begin{center}
	\bgroup
	\def\arraystretch{1.5}
	\begin{tabular}{|p{2.5cm}|p{2cm}|p{8cm}|p{8cm}|p{5cm}|}
		\hline
		\rowcolor{gris}Référence & Exigence & Test & Résultat attendu & Résultat/Commentaires\\
		\hline
		PfcoC.1 & 3.1.1.5.4.1 & Faire un upload, et faire la commande "netstat -a -p" sur le terminal du client. & Le terminal ne doit plus afficher une connexion avec une adresse IP correspondant au serveur.&\vspace*{-0.2cm} \begin{itemize}[label=$\ast$] \item \textcolor{false}{Non implémenté.} \end{itemize} \\
		\hline
	\end{tabular}
	\egroup
\end{center}

\subsection{Download et Upload}

\begin{center}
	\bgroup
	\def\arraystretch{1.5}
	\begin{tabular}{|p{2.5cm}|p{2cm}|p{8cm}|p{8cm}|p{5cm}|}
		\hline
		\rowcolor{gris}Référence & Exigence & Test & Résultat attendu & Résultat/Commentaires\\
		\hline
		DwupC.1 & 3.1.1.6.1 & Connecter le client à un serveur FTP autre qu'un serveur KNS. & La connexion doit s'effectuer. &\vspace*{-0.2cm} \begin{itemize}[label=$\ast$] \item \textcolor{false}{Non impémenté}\end{itemize} \\
		\cline{1-2}\ccline{3-3}{gris}\cline{4-5}
		DwupC.1.1 & 3.1.1.6.3 & Uploader un fichier. & Le fichier est disponible sur le serveur. &\vspace*{-0.2cm} \begin{itemize}[label=$\ast$] \item \textcolor{false}{Non impémenté}\end{itemize}
\\
		\hline
		DwupC.2 & 3.1.1.6.2 & Connecter le client à un serveur FTP KNS avec un compte invalide & La connexion ne doit pas s'effectuer. &\vspace*{-0.2cm} \begin{itemize}[label=$\ast$] \item \textcolor{true}{Validé.}\end{itemize} \\
		\hline
		DwupC.3 & 3.1.1.6.2 & Connecter le client à un serveur FTP KNS avec un compte valide & La connexion doit s'effectuer. &\vspace*{-0.2cm} \begin{itemize}[label=$\ast$] \item \textcolor{true}{Validé.}\end{itemize} \\
		\cline{1-2}\ccline{3-3}{gris}\cline{4-5}
		DwupC.3.1 & 3.1.1.6.4 3.1.1.6.6 & Uploader un fichier, dont l'extension n'est pas interdite. & Une barre de progression s'affiche, arrivée au bout le fichier est disponible sur le serveur. &\vspace*{-0.2cm} \begin{itemize}[label=$\ast$] \item \textcolor{false}{Non impémenté}\end{itemize}
 \\
		\cline{1-2}\ccline{3-3}{gris}\cline{4-5}
		DwupC.3.2 & 3.1.1.6.5 & Uploader un fichier, dont l'extension est interdite. & Un boite de dialogue apparait avec le mention "Le transfert du \textless fichier\_cible\_a\_télécharger \textgreater n'est pas autorisé par le serveur \textless nom\_du\_serveur \textgreater". &\vspace*{-0.2cm} \begin{itemize}[label=$\ast$] \item \textcolor{false}{Non impémenté}\end{itemize}
 \\
		\hline
		DwupC.4 & 3.1.1.6.7 & Uploader ou downloader un fichier. & Une boite de dialogue demande si le transfert doit être crypté. & \vspace*{-0.2cm} \begin{itemize}[label=$\ast$] \item \textcolor{false}{Non impémenté}\end{itemize}
\\		
		\hline	
	\end{tabular}
	\egroup
\end{center}

\subsection{Module de statistique}

\begin{center}
	\bgroup
	\def\arraystretch{1.5}
	\begin{tabular}{|p{2.5cm}|p{2cm}|p{8cm}|p{8cm}|p{5cm}|}
		\hline
		\rowcolor{gris}Référence & Exigence & Test & Résultat attendu & Résultat/Commentaires\\
		\hline
		MdstC.1 & 3.1.3.3.1.1 & Ouvrir le logiciel et se connecter. & Un bouton "Statistiques" est présent dans la barre d'outil. &\vspace*{-0.2cm} \begin{itemize}[label=$\ast$] \item \textcolor{false}{Non impémenté}\end{itemize}\\
		\cline{1-2}\ccline{3-3}{gris}\cline{4-5}
		MdstC.1.1 & 3.1.1.7.1.2 3.1.1.7.2.1 3.1.1.7.2.2 3.1.1.7.2.3 3.1.3.3.1.2 3.1.3.3.1.3 & Cliquer sur le bouton statistique. & Une nouvelle fenêtre apparait avec une partie textuelle dans la moitié haute et une partie graphique dans la moitié basse. &\vspace*{-0.2cm} \begin{itemize}[label=$\ast$] \item \textcolor{false}{Non impémenté}\end{itemize}
\\
		\cline{1-2}\cline{4-5}
		MdstC.1.1.1 & 3.1.1.7.1.2.1 3.1.1.7.2.4 & & Le nom du fichier de l'utilisateur le plus téléchargé apparait dans la partie textuelle. &\vspace*{-0.2cm} \begin{itemize}[label=$\ast$] \item \textcolor{false}{Non impémenté}\end{itemize}
\\
		\cline{1-2}\cline{4-5}
		MdstC.1.1.2 & 3.1.1.7.1.2.2 3.1.1.7.2.5 & & L'identifiant de la personne qui a téléchargé le plus de fichier de l'utilisateur apparait dans la partie textuelle. &\vspace*{-0.2cm} \begin{itemize}[label=$\ast$] \item \textcolor{false}{Non impémenté}\end{itemize}
\\
		\cline{1-2}\cline{4-5}
		MdstC.1.1.3 & 3.1.1.7.1.1 3.1.1.7.1.2.3 3.1.1.7.2.6 & & Un histogramme représentant le nombre de téléchargement de chaque fichier apparait dans la partie graphique. &\vspace*{-0.2cm} \begin{itemize}[label=$\ast$] \item \textcolor{false}{Non impémenté}\end{itemize}
\\
		\cline{1-2}\cline{4-5}
		MdstC.1.1.3.1 & 3.1.1.7.2.7 3.1.1.7.2.8 3.1.1.7.2.9 & & Une barre s'affiche dans l'histogramme pour chaque fichier de l'utilisateur. &\vspace*{-0.2cm} \begin{itemize}[label=$\ast$] \item \textcolor{false}{Non impémenté}\end{itemize}
\\
		\cline{1-2}\cline{4-5}
		MdstC.1.1.3.2 & 3.1.1.7.2.10 & & Chaque barre possède une couleur différente. &\vspace*{-0.2cm} \begin{itemize}[label=$\ast$] \item \textcolor{false}{Non impémenté}\end{itemize}
\\
		\cline{1-2}\cline{4-5}
		MdstC.1.1.3.3 & 3.1.1.7.2.11 & & La hauteur des barres correspond au nombre de téléchargement de chaque fichier. &\vspace*{-0.2cm} \begin{itemize}[label=$\ast$] \item \textcolor{false}{Non impémenté}\end{itemize}
\\
		\cline{1-2}\cline{4-5}
		MdstC.1.1.3.4 & 3.1.1.7.2.12 & & La largeur des barres est adaptée pour que l'histogramme prenne entièrement la largeur de la page. &\vspace*{-0.2cm} \begin{itemize}[label=$\ast$] \item \textcolor{false}{Non impémenté}\end{itemize}
\\
		\cline{1-2}\cline{4-5}
		MdstC.1.1.3.5 & 3.1.1.7.2.13 & & En ordonnée le nombre de fichier est affiché. &\vspace*{-0.2cm} \begin{itemize}[label=$\ast$] \item \textcolor{false}{Non impémenté}\end{itemize}
\\
		\cline{1-2}\cline{4-5}
		MdstC.1.1.3.6 & 3.1.1.7.2.14 & & En abscisse un nombre est affiché sous chacune des barres. &\vspace*{-0.2cm} \begin{itemize}[label=$\ast$] \item \textcolor{false}{Non impémenté}\end{itemize}
\\
		\cline{1-2}\cline{4-5}
		MdstC.1.1.3.6.1 & 3.1.1.7.2.15 & & Une légende est affichée avec la correspondance numéro de barre/nom du fichier. &\vspace*{-0.2cm} \begin{itemize}[label=$\ast$] \item \textcolor{false}{Non impémenté}\end{itemize}
\\
		\hline
	\end{tabular}
	\egroup
\end{center}
\newpage
\begin{center}
	\bgroup
	\def\arraystretch{1.5}
	\begin{tabular}{|p{2.5cm}|p{2cm}|p{8cm}|p{8cm}|p{5cm}|}
		\hline
		\rowcolor{gris}Référence & Exigence & Test & Résultat attendu & Résultat/Commentaires\\
		\hline
		\cline{1-2}\cline{4-5}
		MdstC.1.2 & 3.1.3.3.1.4 & & Le titre "Statistiques" apparait en haut de la fenêtre. &\vspace*{-0.2cm} \begin{itemize}[label=$\ast$] \item \textcolor{false}{Non impémenté}\end{itemize}
\\
		\cline{1-2}\cline{4-5}
		MdstC.1.3 & 3.1.3.3.1.5 3.1.3.3.1.6 3.1.3.3.1.8 & & Une barre d'outil apparait en haut de la fenêtre avec un bouton "Exporter" à gauche. Un bouton "Fermer" apparait en bas à droite. &\vspace*{-0.2cm} \begin{itemize}[label=$\ast$] \item \textcolor{false}{Non impémenté}\end{itemize}
\\
		\cline{1-2}\ccline{3-3}{gris}\cline{4-5}
		MdstC.1.3.1 & 3.1.3.3.1.7 & Cliquer sur le bouton "Exporter". & Un fichier html est crée et contient le rapport des statistiques de l'utilisateur. &\vspace*{-0.2cm} \begin{itemize}[label=$\ast$] \item \textcolor{false}{Non impémenté}\end{itemize}
\\
		\cline{1-2}\ccline{3-3}{gris}\cline{4-5}
		MdstC.1.3.2 & 3.1.3.3.1.8 & Cliquer sur le bouton "Fermer". & La fenêtre se ferme. &\vspace*{-0.2cm} \begin{itemize}[label=$\ast$] \item \textcolor{false}{Non impémenté}\end{itemize}
\\
		\hline
	\end{tabular}
	\egroup
\end{center}

\subsection{Cryptage}

\begin{center}
	\bgroup
	\def\arraystretch{1.5}
	\begin{tabular}{|p{2.5cm}|p{2cm}|p{8cm}|p{8cm}|p{5cm}|}
		\hline
		\rowcolor{gris}Référence & Exigence & Test & Résultat attendu & Résultat/Commentaires\\
		\hline
		CrypC.1 & 3.1.3.1.2 & Faire un clic droit sur un fichier. & Le menu de gestion de fichier apparait avec le bouton "Cryptage du fichier" à l'intérieur. &\vspace*{-0.2cm} \begin{itemize}[label=$\ast$] \item \textcolor{false}{Non impémenté}\end{itemize}\\
		\cline{1-2}\ccline{3-3}{gris}\cline{4-5}
		CrypC.1.1 & 3.1.3.1.1 3.1.3.1.3 3.1.3.1.4 & Cliquer sur le bouton "Cryptage du fichier". & Le fichier est crypté et illisible pour les utilisateurs ne connaissant pas la clé de cryptage. &\vspace*{-0.2cm} \begin{itemize}[label=$\ast$] \item \textcolor{false}{Non impémenté}\end{itemize}
\\
		\hline
		CrypC.2 & 3.1.3.1.5 3.1.3.1.6 & Télécharger un fichier crypté dont la clé de cryptage est enregistrée. & Le fichier décrypté est téléchargé. &\vspace*{-0.2cm} \begin{itemize}[label=$\ast$] \item \textcolor{false}{Non impémenté}\end{itemize}
\\
		\hline
		CrypC.3 & 3.1.3.1.7 & Télécharger un fichier crypté dont la clé de cryptage n'est pas enregistrée. & Une boite de dialogue s'affiche avec un champ de saisie "Clé". &\vspace*{-0.2cm} \begin{itemize}[label=$\ast$] \item \textcolor{false}{Non impémenté}\end{itemize}
\\
		\cline{1-2}\ccline{3-3}{gris}\cline{4-5}
		CrypC.3.1 & 3.1.3.1.8 & Entrer la clé de cryptage dans le champ. & Le fichier décrypté est téléchargé. &\vspace*{-0.2cm} \begin{itemize}[label=$\ast$] \item \textcolor{false}{Non impémenté}\end{itemize}
\\
		\cline{1-2}\ccline{3-3}{gris}\cline{4-5}
		CrypC.3.1.1 & 3.1.3.1.9 & Recommencer le téléchargement. & Le fichier décrypté est téléchargé sans redemander la clé de cryptage. &\vspace*{-0.2cm} \begin{itemize}[label=$\ast$] \item \textcolor{false}{Non impémenté}\end{itemize}
\\
		\hline
		CrypC.4 & 3.1.3.1.10 & Crypter un fichier sur le serveur et le télécharger sur un client. & Le client entre une clé, si celle ci est en accord avec celle sur le serveur, le fichier est décrypté. &\vspace*{-0.2cm} \begin{itemize}[label=$\ast$] \item \textcolor{false}{Non impémenté}\end{itemize}
\\
		\hline
		CrypC.5 & 3.1.3.1.11 & Ouvrir le fichier contenant les clés. & Les clés ne sont pas écrit en clair. &\vspace*{-0.2cm} \begin{itemize}[label=$\ast$] \item \textcolor{false}{Non impémenté}\end{itemize}
\\
		\hline	
	\end{tabular}
	\egroup
\end{center}







\section{Tests Serveur}
\subsection{Authentification}
\begin{center}
	\bgroup
	\def\arraystretch{1.5}
	\begin{tabular}{|p{2.5cm}|p{2cm}|p{8cm}|p{8cm}|p{5cm}|}
		\hline
		\rowcolor{gris}Référence & Exigence & Test & Résultat attendu & Résultat/Commentaires\\
		\cline{1-2}\ccline{3-3}{gris}\cline{4-5}
		AuthS.1 & 3.1.2.1.2.1 & Envoyer une chaine de caractère contenant un pseudo suivit d'un mot de passe & Le programme détecte cette arrivée et l'indique dans une console de log &\vspace*{-0.2cm} \begin{itemize}[label=$\ast$] \item \textcolor{true}{Validé.}\end{itemize} \\
		\cline{1-2}\ccline{3-3}{gris}\cline{4-5}		
		AuthS.2 & 3.1.2.1.2.2 & Après avoir récupéré les informations essayer de les faire correspondre à la liste des pseudo / Mot de passe. & Le programme indique dans une console de log le succès ou l'échec de l'opération &\vspace*{-0.2cm} \begin{itemize}[label=$\ast$] \item \textcolor{true}{Validé.}\end{itemize} \\
		\cline{1-2}\ccline{3-3}{gris}\cline{4-5}		
		AuthS.3 & 3.1.2.1.2.3 & Envoyer une chaine contenant le mot "succès" ou "échec" (suivant la réponse au test AuthS.2) sur un client sommaire. & Chaine reçu en concordance avec le résultat du test AuthS.2 &\vspace*{-0.2cm} \begin{itemize}[label=$\ast$] \item \textcolor{true}{Validé.}\end{itemize} \\
		\cline{1-2}\ccline{3-3}{gris}\cline{4-5}		
		AuthS.4 & 3.1.2.1.2.4 & Vérifier que le pseudo et le mot de passe admin sont unique. & Les identifiants sont unique. &\vspace*{-0.2cm} \begin{itemize}[label=$\ast$] \item \textcolor{false}{Non impémenté}\end{itemize}
 \\
		\cline{1-2}\ccline{3-3}{gris}\cline{4-5}
		AuthS.5 & 3.1.2.1.2.5 & Se connecter en temps qu'Administrateur. &  Avoir accès à la totalité des fichiers du serveur. &\vspace*{-0.2cm} \begin{itemize}[label=$\ast$] \item \textcolor{false}{Non impémenté}\end{itemize} \\
		\hline
	\end{tabular}
	\egroup
\end{center}

\subsection{Partage de fichier}
\begin{center}
	\bgroup
	\def\arraystretch{1.5}
	\begin{tabular}{|p{2.5cm}|p{2cm}|p{8cm}|p{8cm}|p{5cm}|}
		\hline
		\rowcolor{gris}Référence & Exigence & Test & Résultat attendu & Résultat/Commentaires\\
		\hline

		PafiS.1 & 3.1.2.2.1 & Créer un compte et essayer de partager des fichiers présents à partir de la racine. & Les fichiers ne doivent pas être accessible. &\vspace*{-0.2cm} \begin{itemize}[label=$\ast$] \item \textcolor{false}{non implémenté.}\end{itemize} \\
		\cline{1-2}\ccline{3-3}{gris}\cline{4-5}
		PafiS.2 & 3.1.2.2.2 & Crée 2 comptes et partager un même fichier à ses 2 comptes & Se connecter sur les 2 comptes et vérifier l'accès aux fichiers. &\vspace*{-0.2cm} \begin{itemize}[label=$\ast$] \item \textcolor{false}{non implémenté.}\end{itemize}
\\
		\cline{1-2}\ccline{3-3}{gris}\cline{4-5}
		PafiS.3 & 3.1.2.2.3 & Crée un compte lui partager un fichier et bloqué l'accès à un sous dossier & Se connecter sur le compte crée, le fichier ne doit pas être présent dans la liste des fichiers. &\vspace*{-0.2cm} \begin{itemize}[label=$\ast$] \item \textcolor{false}{non implémenté.}\end{itemize}
\\
		\cline{1-2}\ccline{3-3}{gris}\cline{4-5}
		PafiS.4 & 3.1.2.2.4 & Ouvrir un explorateur de fichier sur le serveur & Les fichiers sont visible avec l'explorateur de fichier &\vspace*{-0.2cm} \begin{itemize}[label=$\ast$] \item \textcolor{false}{non implémenté.}\end{itemize}
\\
		\hline
	\end{tabular}
	\egroup
\end{center}

\subsection{Gestionnaire de comptes utilisateurs}
\begin{center}
	\bgroup
	\def\arraystretch{1.5}
	\begin{tabular}{|p{2.5cm}|p{2cm}|p{8cm}|p{8cm}|p{5cm}|}
		\hline
		\rowcolor{gris}Référence & Exigence & Test & Résultat attendu & Résultat/Commentaires\\
		\hline
		GecuS.1 & 3.1.2.3.1.1 & Ouvrir le logiciel & Une zone est présente à gauche. & \vspace*{-0.2cm} \begin{itemize}[label=$\ast$] \item \textcolor{false}{Non impémenté}\end{itemize}
\\
		\cline{1-2}\cline{4-5}
		GecuS.1.1 & 3.1.2.3.1.2 & & Dans cette zone se trouve la liste des comptes. &\vspace*{-0.2cm} \begin{itemize}[label=$\ast$] \item \textcolor{false}{Non impémenté}\end{itemize}
\\
		\cline{1-2}\cline{4-5}
		GecuS.1.2 & 3.1.2.3.1.3 & & En face de chaque compte se trouve un voyant vert si la personne est connectée, rouge sinon. &\vspace*{-0.2cm} \begin{itemize}[label=$\ast$] \item \textcolor{false}{Non impémenté}\end{itemize}
\\
		\cline{1-2}\cline{4-5}
		GecuS.1.3 & 3.1.2.3.1.4 & & En bas de la zone de gauche 2 boutons sont présent &\vspace*{-0.2cm} \begin{itemize}[label=$\ast$] \item \textcolor{false}{Non impémenté}\end{itemize}
\\
		\cline{1-2}\ccline{3-3}{gris}\cline{4-5}
		GecuS.1.3.1 & 3.1.2.3.1.4.1 & & Un bouton "-" rouge est présent en bas à droite &\vspace*{-0.2cm} \begin{itemize}[label=$\ast$] \item \textcolor{false}{Non impémenté}\end{itemize}
\\
		\cline{1-2}\cline{4-5}
		GecuS.1.3.1.1 & 3.1.2.3.1.4.2 & Cliquer sur le bouton "-" & Un pop-up apparait. &\vspace*{-0.2cm} \begin{itemize}[label=$\ast$] \item \textcolor{false}{Non impémenté}\end{itemize}
\\
		\cline{1-2}\cline{4-5}
		GecuS.1.3.2 & 3.1.2.3.1.4.3 & & Un bouton "+" vert est présent en bas à droite &\vspace*{-0.2cm} \begin{itemize}[label=$\ast$] \item \textcolor{false}{Non impémenté}\end{itemize}
\\
		\cline{1-2}\ccline{3-3}{gris}\cline{4-5}
		GecuS.1.3.2.1 & 3.1.2.3.1.4.4 & Cliquer sur le bouton "+" & Un pop-up apparait. &\vspace*{-0.2cm} \begin{itemize}[label=$\ast$] \item \textcolor{false}{Non impémenté}\end{itemize}
\\
		\cline{1-2}\cline{4-5}
		GecuS.1.3.2.1.1 & 3.1.2.6.2.1 & & Le pop-up contient 4 champs &\vspace*{-0.2cm} \begin{itemize}[label=$\ast$] \item \textcolor{false}{Non impémenté}\end{itemize}
\\
		\cline{1-2}\cline{4-5}
		GecuS.1.3.2.1.1.1 & 3.2.6.2.1.1 & & Un champ texte aligné au texte "Nom d'utilisateur" &\vspace*{-0.2cm} \begin{itemize}[label=$\ast$] \item \textcolor{false}{Non impémenté}\end{itemize}
\\
		\cline{1-2}\cline{4-5}
		GecuS.1.3.2.1.1.2 & 3.2.6.2.1.2 & & Un champ texte aligné au texte "Mot de passe" &\vspace*{-0.2cm} \begin{itemize}[label=$\ast$] \item \textcolor{false}{Non impémenté}\end{itemize}
\\
		\cline{1-2}\cline{4-5}
		GecuS.1.3.2.1.1.3 & 3.2.6.2.1.3 & & Un champs de type sélection de fichier &\vspace*{-0.2cm} \begin{itemize}[label=$\ast$] \item \textcolor{false}{Non impémenté}\end{itemize}
\\
		\cline{1-2}\cline{4-5}
		GecuS.1.3.2.1.1.3 & 3.2.6.2.1.4 & & Une case à cochée aligné avec le texte "Utilisateur temporaire" &\vspace*{-0.2cm} \begin{itemize}[label=$\ast$] \item \textcolor{false}{Non impémenté}\end{itemize}
\\
		\cline{1-2}\ccline{3-3}{gris}\cline{4-5}
		GecuS.1.3.2.1.2 & 3.2.6.2.2 & Cocher la case "Utilisateur Temporaire" & Un champ de type date s'ajoute" &\vspace*{-0.2cm} \begin{itemize}[label=$\ast$] \item \textcolor{false}{Non impémenté}\end{itemize}
\\
		\cline{1-2}\cline{4-5}
		GecuS.1.3.3 & 3.2.6.2.3 & & En bas à droite se trouve un bouton annulé &\vspace*{-0.2cm} \begin{itemize}[label=$\ast$] \item \textcolor{false}{Non impémenté}\end{itemize}
\\
		\cline{1-2}\cline{4-5}
		GecuS.1.3.4 & 3.2.6.2.4 & & En bas à gauche un bouton "Valider" est présent &\vspace*{-0.2cm} \begin{itemize}[label=$\ast$] \item \textcolor{false}{Non impémenté}\end{itemize}
\\ 
		\cline{1-2}\cline{4-5}
		GecuS.1.4 & 3.1.2.3.1.5 & & Les comptes temporaires sont soulignés &\vspace*{-0.2cm} \begin{itemize}[label=$\ast$] \item \textcolor{false}{Non impémenté}\end{itemize}
\\
		\cline{1-2}\cline{4-5}
		GecuS.1.5 & 3.1.2.3.1.6 & & Le compte administrateur n'est pas présent dans la liste &\vspace*{-0.2cm} \begin{itemize}[label=$\ast$] \item \textcolor{false}{Non impémenté}\end{itemize}
\\
		\cline{1-2}\cline{4-5}
		GecuS.1.6 & 3.1.2.3.1.7 & & Les comptes bannis sont en rouge &\vspace*{-0.2cm} \begin{itemize}[label=$\ast$] \item \textcolor{false}{Non impémenté}\end{itemize}
\\
		\cline{1-2}\cline{4-5}
		GecuS.1.7 & 3.1.2.3.1.8 & & Les comptes bannis temporairement sont en gris clair &\vspace*{-0.2cm} \begin{itemize}[label=$\ast$] \item \textcolor{false}{Non impémenté}\end{itemize}
\\
		\cline{1-2}\cline{4-5}
		GecuS.1.8 & 3.1.2.3.1.9 & & Les comptes bannis partiellement sont en orange. &\vspace*{-0.2cm} \begin{itemize}[label=$\ast$] \item \textcolor{false}{Non impémenté}\end{itemize}
\\
		\cline{1-2}\cline{4-5}
		GecuS.1.9 & 3.1.2.3.1.10 & & Les comptes normaux sont en noir &\vspace*{-0.2cm} \begin{itemize}[label=$\ast$] \item \textcolor{false}{Non impémenté}\end{itemize}
\\
		\cline{1-2}\ccline{3-3}{gris}\cline{4-5}
		GecuS.1.10 & 3.1.2.3.1.11 & Cliquer sur un compte. & La liste de droite change et contient les répertoire auquel le compte à accès. &\vspace*{-0.2cm} \begin{itemize}[label=$\ast$] \item \textcolor{false}{Non impémenté}\end{itemize}
\\
		\hline
	\end{tabular}
	\egroup
\end{center}
\newpage

\begin{center}
	\bgroup
	\def\arraystretch{1.5}
	\begin{tabular}{|p{2.5cm}|p{2cm}|p{8cm}|p{8cm}|p{5cm}|}
		\hline
		\rowcolor{gris}Référence & Exigence & Test & Résultat attendu & Résultat/Commentaires\\
		\hline
		GecuS.2 & 3.1.2.3.1 & Se connecter en tant qu'administrateur et essayer de créer un compte & Le compte peut être crée et les droits donné. &\vspace*{-0.2cm} \begin{itemize}[label=$\ast$] \item \textcolor{false}{Non impémenté}\end{itemize}
\\
		\cline{1-2}\ccline{3-3}{gris}\cline{4-5}
		GecuS.3 & 3.1.2.3.2 & Se connecter en tant qu'administrateur et crée un compte et cocher compte temporaire & Un compte temporaire est crée &\vspace*{-0.2cm} \begin{itemize}[label=$\ast$] \item \textcolor{false}{Non impémenté}\end{itemize}
 \\
		\cline{1-2}\ccline{3-3}{gris}\cline{4-5}
		GecuS.4 & 3.1.2.3.3 & Regarder la liste des membres présente à droite & Elle contient la totalité des comptes crée &\vspace*{-0.2cm} \begin{itemize}[label=$\ast$] \item \textcolor{false}{Non impémenté}\end{itemize}
 \\
		\cline{1-2}\ccline{3-3}{gris}\cline{4-5}
		GecuS.5 & 3.1.2.3.4 & Cliquer sur un compte & La liste des fichiers auquel il a accès se trouve à droite et un plus est présent en bas à droite. &\vspace*{-0.2cm} \begin{itemize}[label=$\ast$] \item \textcolor{false}{Non impémenté}\end{itemize}
 \\
		\cline{1-2}\ccline{3-3}{gris}\cline{4-5}
		GecuS.6 & 3.1.2.3.5 & Cliquer sur un compte et cliquer sur le bouton "-" présent en bas à droite de la partie droite & Le compte est banni après confirmation. &\vspace*{-0.2cm} \begin{itemize}[label=$\ast$] \item \textcolor{false}{Non impémenté}\end{itemize}
 \\
		\cline{1-2}\cline{4-5}
		GecuS.7 & 3.1.2.3.6 & & Le compte est supprimé après confirmation. &\vspace*{-0.2cm} \begin{itemize}[label=$\ast$] \item \textcolor{false}{Non impémenté}\end{itemize}
 \\
		\hline
	\end{tabular}
	\egroup
\end{center}

\subsubsection{Bannissement d’un compte}
\begin{center}
	\bgroup
	\def\arraystretch{1.5}
	\begin{tabular}{|p{2.5cm}|p{2cm}|p{8cm}|p{8cm}|p{5cm}|}
		\hline
		\rowcolor{gris}Référence & Exigence & Test & Résultat attendu & Résultat/Commentaires\\
		\hline
		BancS.1 & 3.1.2.4.1 & Cliquer sur un compte. & Un moins est présent dans la zone de droites. &\vspace*{-0.2cm} \begin{itemize}[label=$\ast$] \item \textcolor{false}{Non impémenté}\end{itemize}\\
		\cline{1-2}\ccline{3-3}{gris}\cline{4-5}
		BancS.1.1 & 3.1.2.6.1.1 & Cliquer sur le moins & Un pop-up s'ouvre, celui ci contient 2 cases &\vspace*{-0.2cm} \begin{itemize}[label=$\ast$] \item \textcolor{false}{Non impémenté}\end{itemize}
\\
		\cline{1-2}\cline{4-5}
		BancS.1.1.1 & 3.1.2.6.1.1.1 & & Une case à cocher alignée avec le texte "Supprimer le compte" &\vspace*{-0.2cm} \begin{itemize}[label=$\ast$] \item \textcolor{false}{Non impémenté}\end{itemize}
\\
		\cline{1-2}\cline{4-5}
		BancS.1.1.2 & 3.1.2.6.1.1.2 & & Une case à cocher alignée avec le texte "Bannir le compte" &\vspace*{-0.2cm} \begin{itemize}[label=$\ast$] \item \textcolor{false}{Non impémenté}\end{itemize}
\\
		\cline{1-2}\ccline{3-3}{gris}\cline{4-5}
		BancS.1.2 & 3.1.2.6.1.2 & Cocher une case puis l'autre & Une seul case doit être cochée &\vspace*{-0.2cm} \begin{itemize}[label=$\ast$] \item \textcolor{false}{Non impémenté}\end{itemize}
\\
		\cline{1-2}\cline{4-5}
		BancS.1.3 & 3.1.2.6.1.3 & & Le pop-up contient 2 boutons en bas &\vspace*{-0.2cm} \begin{itemize}[label=$\ast$] \item \textcolor{false}{Non impémenté}\end{itemize}
\\
		\cline{1-2}\cline{4-5}
		BancS.1.3.1 & 3.1.2.6.1.3.1 & & Un bouton "Valider" &\vspace*{-0.2cm} \begin{itemize}[label=$\ast$] \item \textcolor{false}{Non impémenté}\end{itemize}
\\
		\cline{1-2}\cline{4-5}
		BancS.1.3.2 & 3.1.2.6.1.3.2 & & Un bouton "Annuler" &\vspace*{-0.2cm} \begin{itemize}[label=$\ast$] \item \textcolor{false}{Non impémenté}\end{itemize}
\\
		\cline{1-2}\ccline{3-3}{gris}\cline{4-5}
		BancS.1.4 & 3.1.2.6.1.4 & Cocher "Supprimer le compte" & Seul les 2 cases ( dont "supprimer le compte" cochée ) plus les 2 boutons sont présent. &\vspace*{-0.2cm} \begin{itemize}[label=$\ast$] \item \textcolor{false}{Non impémenté}\end{itemize}
\\
		\cline{1-2}\ccline{3-3}{gris}\cline{4-5}
		BancS.1.5 & 3.1.2.6.1.5 & Cocher "Bannir le compte" & 2 Champs sont ajouté &\vspace*{-0.2cm} \begin{itemize}[label=$\ast$] \item \textcolor{false}{Non impémenté}\end{itemize}
\\
		\cline{1-2}\cline{4-5}
		BancS.1.5.1 & 3.1.2.6.1.5.1 & & Un champ Motif est ajouté aligné au texte "Motif de bannissement" &\vspace*{-0.2cm} \begin{itemize}[label=$\ast$] \item \textcolor{false}{Non impémenté}\end{itemize}
\\
		\cline{1-2}\cline{4-5}
		BancS.1.5.2 & 3.1.2.6.1.5.2 & & Un case à cocher est alignée avec le texte "Bannissement temporaire" &\vspace*{-0.2cm} \begin{itemize}[label=$\ast$] \item \textcolor{false}{Non impémenté}\end{itemize}
\\
		\cline{1-2}\ccline{3-3}{gris}\cline{4-5}
		BancS.1.5.2.1 & 3.1.2.6.1.6 & Cocher la case "Bannissement temporaire" & Un champs de type date s'ajoute en dessous de la case. &\vspace*{-0.2cm} \begin{itemize}[label=$\ast$] \item \textcolor{false}{Non impémenté}\end{itemize}
\\ 
		\cline{1-2}\ccline{3-3}{gris}\cline{4-5}
		BancS.1.6 & 3.1.2.6.1.7 & Ouvrir le fichier .accban & Il contient le motif de bannissement &\vspace*{-0.2cm} \begin{itemize}[label=$\ast$] \item \textcolor{false}{Non impémenté}\end{itemize}
\\
		\cline{1-2}\cline{4-5}
		BancS.1.6.1 & 3.1.2.6.1.8 & & Il contient une durée en seconde &\vspace*{-0.2cm} \begin{itemize}[label=$\ast$] \item \textcolor{false}{Non impémenté}\end{itemize}
\\
		\cline{1-2}\ccline{3-3}{gris}\cline{4-5}
		BancS.2 & 3.1.2.4.2 & Et retirer les droits sur un fichier. & Ouvrir le fichier .acc et vérifier que les droits sur le fichier ont été retiré. &\vspace*{-0.2cm} \begin{itemize}[label=$\ast$] \item \textcolor{false}{Non impémenté}\end{itemize}
\\
		\cline{1-2}\ccline{3-3}{gris}\cline{4-5}
		BancS.3 & 3.1.2.4.3 & Bannir temporairement une personne. & Vérifier que le fichier .accban est mis à jours. &\vspace*{-0.2cm} \begin{itemize}[label=$\ast$] \item \textcolor{false}{Non impémenté}\end{itemize}
\\
		\cline{1-2}\ccline{3-3}{gris}\cline{4-5}
		BancS.4 & 3.1.2.4.4 & Bannir un compte. & Vérifier que le champs raison est présent lors du bannissement. &\vspace*{-0.2cm} \begin{itemize}[label=$\ast$] \item \textcolor{false}{Non impémenté}\end{itemize}
\\
		\cline{1-2}\ccline{3-3}{gris}\cline{4-5}
		BancS.5 & 3.1.2.4.5 & Supprimer les droits d'un compte sur un fichier. & Vérifier que l'utilisateur n'a plus les droits sur le fichier. &\vspace*{-0.2cm} \begin{itemize}[label=$\ast$] \item \textcolor{false}{Non impémenté}\end{itemize}
\\
		\hline
	\end{tabular}
	\egroup
\end{center}

\subsubsection{Bannissement temporaire}
\begin{center}
	\bgroup
	\def\arraystretch{1.5}
	\begin{tabular}{|p{2.5cm}|p{2cm}|p{8cm}|p{8cm}|p{5cm}|}
		\hline
		\rowcolor{gris}Référence & Exigence & Test & Résultat attendu & Résultat/Commentaires\\
		\hline
		BantS.1 & 3.1.2.5.1 & Cliquer sur un compte puis sur "-" présent sur la zone de gauche. & Un champ permet de donner une durée. &\vspace*{-0.2cm} \begin{itemize}[label=$\ast$] \item \textcolor{false}{Non impémenté}\end{itemize}\\
		\cline{1-2}\ccline{3-3}{gris}\cline{4-5}
		BantS.2 & 3.1.2.5.2 & Cliquer sur un compte bannis. & Le temps restant apparait à droite. &\vspace*{-0.2cm} \begin{itemize}[label=$\ast$] \item \textcolor{false}{Non impémenté}\end{itemize}
\\
		\cline{1-2}\ccline{3-3}{gris}\cline{4-5}
		BantS.3 & 3.1.2.5.3 & Cliquer sur un compte puis sur "-" présent dans la zone de gauche. & Un champs permet de donner la raison. &\vspace*{-0.2cm} \begin{itemize}[label=$\ast$] \item \textcolor{false}{Non impémenté}\end{itemize}
\\
		\hline
	\end{tabular}
	\egroup
\end{center}

\subsubsection{Bannissement partiel}
\begin{center}
	\bgroup
	\def\arraystretch{1.5}
	\begin{tabular}{|p{2.5cm}|p{2cm}|p{8cm}|p{8cm}|p{5cm}|}
		\hline
		\rowcolor{gris}Référence & Exigence & Test & Résultat attendu & Résultat/Commentaires\\
		\hline
		BanpS.1 & 3.1.2.6.3.1 & Ouvrir le logiciel & A droite se trouve la liste de tout les fichiers présent sur le serveur &\vspace*{-0.2cm} \begin{itemize}[label=$\ast$] \item \textcolor{false}{Non impémenté}\end{itemize}\\
		\cline{1-2}\ccline{3-3}{gris}\cline{4-5}
		BanpS.2 & 3.1.2.6.3.2 & Cliquer sur un utilisateur & Dans la partie droite seul la liste des fichiers accessible par l'utilisateur est présente. &\vspace*{-0.2cm} \begin{itemize}[label=$\ast$] \item \textcolor{false}{Non impémenté}\end{itemize}
\\
		\cline{1-2}\cline{4-5}
		BanpS.2.1 & 3.1.2.6.3.3 & & 2 Champs sont présent en bas à droite de la partie de droite. &\vspace*{-0.2cm} \begin{itemize}[label=$\ast$] \item \textcolor{false}{Non impémenté}\end{itemize}
\\
		\cline{1-2}\ccline{3-3}{gris}\cline{4-5}
		BanpS.2.1.1 & 3.1.2.6.3.3.1 & & Un "+" vert est présent en bas à droite &\vspace*{-0.2cm} \begin{itemize}[label=$\ast$] \item \textcolor{false}{Non impémenté}\end{itemize}
\\
		\cline{1-2}\cline{4-5}
		BanpS.2.1.1.1 & 3.1.2.6.3.5 & Cliquer sur le "+" & Un pop-up apparait, celui ci permettant de sélectionner des fichiers &\vspace*{-0.2cm} \begin{itemize}[label=$\ast$] \item \textcolor{false}{Non impémenté}\end{itemize}
\\
		\cline{1-2}\ccline{3-3}{gris}\cline{4-5}
		BanpS.2.1.2 & 3.1.2.6.3.3.1 & & A droite du plus un "-" Rouge est présent &\vspace*{-0.2cm} \begin{itemize}[label=$\ast$] \item \textcolor{false}{Non impémenté}\end{itemize}
\\
		\cline{1-2}\ccline{3-3}{gris}\cline{4-5}
		BanpS.3 & 3.1.2.6.1 & Cliquer sur un compte puis sélectionner un fichier dans la zone de droite et cliquer sur "-" & L'utilisateur n'a plus les droits sur le fichier. &\vspace*{-0.2cm} \begin{itemize}[label=$\ast$] \item \textcolor{false}{Non impémenté}\end{itemize}
\\
		\cline{1-2}\ccline{3-3}{gris}\cline{4-5}
		BanpS.4 & 3.1.2.6.2 & Donner une raison au bannissement. & Vérifier que le motif est présent dans le fichier .accban. &\vspace*{-0.2cm} \begin{itemize}[label=$\ast$] \item \textcolor{false}{Non impémenté}\end{itemize}
\\
		\cline{1-2}\ccline{3-3}{gris}\cline{4-5}
		BanpS.5 & 3.1.2.6.3 & Cliquer sur un compte puis Cliquer sur le moins de la partie droite. & Un champs permet de donner la durée du bannissement.  &\vspace*{-0.2cm} \begin{itemize}[label=$\ast$] \item \textcolor{false}{Non impémenté}\end{itemize}
\\
		\hline
	\end{tabular}
	\egroup
\end{center}


\subsection{Stockage des comptes}
\begin{center}
	\bgroup
	\def\arraystretch{1.5}
	\begin{tabular}{|p{2.5cm}|p{2cm}|p{8cm}|p{8cm}|p{5cm}|}
		\hline
		\rowcolor{gris}Référence & Exigence & Test & Résultat attendu & Résultat/Commentaires\\
		\hline
		StcS.1 & 3.1.2.6.4.1 & Ouvrir les fichiers contenant les comptes & Les mots de passe ne sont pas écrit en clair. &\vspace*{-0.2cm} \begin{itemize}[label=$\ast$] \item \textcolor{true}{Validé.}\end{itemize} \\
		\cline{1-2}\cline{4-5}
		StcS.2 & 3.1.2.6.4.2 & & Sur chaque ligne est présent les droits associé à un compte. &\vspace*{-0.2cm} \begin{itemize}[label=$\ast$] \item \textcolor{false}{Non impémenté}\end{itemize} \\
		\cline{1-2}\cline{4-5}
		StcS.3 & 3.1.2.6.4.3 & & Les comptes de differents types sont dans des fichiers séparés. &\vspace*{-0.2cm} \begin{itemize}[label=$\ast$] \item \textcolor{false}{Non impémenté}\end{itemize} \\
		\cline{1-2}\ccline{3-3}{gris}\cline{4-5}
		StcS.4 & 3.1.2.6.4.4 & Se connecter sur un compte non présent dans l'un des 3 fichiers & Le programme nous indique que le mot de passe et / ou le nom de compte est erroné. &\vspace*{-0.2cm} \begin{itemize}[label=$\ast$] \item \textcolor{true}{Validé.}\end{itemize} \\
		\hline
	\end{tabular}
	\egroup
\end{center}

\subsubsection{Compte Administrateur}
\begin{center}
	\bgroup
	\def\arraystretch{1.5}
	\begin{tabular}{|p{2.5cm}|p{2cm}|p{8cm}|p{8cm}|p{5cm}|}
		\hline
		\rowcolor{gris}Référence & Exigence & Test & Résultat attendu & Résultat/Commentaires\\
		\hline
		CoaS.1 & 3.1.2.6.5.1 & Ouvrir le fichier contenant les comptes administrateurs & Vérifier qu'il n'y a qu'une seule ligne dedans. &\vspace*{-0.2cm} \begin{itemize}[label=$\ast$] \item \textcolor{false}{Non impémenté}\end{itemize} \\
		\cline{1-2}\cline{4-5}
		CoaS.2 & 2.1.2.6.5.2 & & Seul le compte administrateur est présent dans le fichier. &\vspace*{-0.2cm} \begin{itemize}[label=$\ast$] \item \textcolor{false}{Non impémenté}\end{itemize}
\\
		\cline{1-2}\ccline{3-3}{gris}\cline{4-5}
		CoaS.3 & 3.1.2.6.5.3 & Ouvrir le dossier contenant les comptes administrateur & Le fichier ne possède pas d'extension. &\vspace*{-0.2cm} \begin{itemize}[label=$\ast$] \item \textcolor{false}{Non impémenté}\end{itemize}
\\
		\cline{1-2}\cline{4-5}
		CoaS.4 & 3.1.2.6.5.4 & & Le fichier s'appel .shadow &\vspace*{-0.2cm} \begin{itemize}[label=$\ast$] \item \textcolor{false}{Non impémenté}\end{itemize}
\\
		\cline{1-2}\ccline{3-3}{gris}\cline{4-5}
		CoaS.5 & 3.1.2.6.5.5 & Ouvrir le fichier contenant les comptes administrateurs & Le fichier contient une ligne unique de la forme : "pseudo@motDePasse:/CheminDeReference". &\vspace*{-0.2cm} \begin{itemize}[label=$\ast$] \item \textcolor{false}{Non impémenté}\end{itemize}
\\
		\cline{1-2}\ccline{3-3}{gris}\cline{4-5}
		CoaS.6 & 3.1.2.6.5.6 & Essayer de donner des droits à partir de la racine. & Impossible d'accéder aux fichier présent avant la racine présente dans le fichier .shadow. &\vspace*{-0.2cm} \begin{itemize}[label=$\ast$] \item \textcolor{false}{Non impémenté}\end{itemize}
\\
		\cline{1-2}\ccline{3-3}{gris}\cline{4-5}
		CoaS.7 & 3.1.2.6.5.7 & Supprimer un fichier en même temps qu'un utilisateur le visite. & L'utilisateur n'a plus d'accès au fichier. &\vspace*{-0.2cm} \begin{itemize}[label=$\ast$] \item \textcolor{false}{Non impémenté}\end{itemize}
\\
		\cline{1-2}\ccline{3-3}{gris}\cline{4-5}
		CoaS.8 & 3.1.2.6.5.8 & Essayer d'accéder au compte administrateur à partir d'un client & Impossibilité de se connecter au compte. &\vspace*{-0.2cm} \begin{itemize}[label=$\ast$] \item \textcolor{false}{Non impémenté}\end{itemize}
\\
		\cline{1-2}\cline{4-5}
		CoaS.9 & 3.1.2.6.5.9 & & Impossibilité de se connecter au compte. &\vspace*{-0.2cm} \begin{itemize}[label=$\ast$] \item \textcolor{false}{Non impémenté}\end{itemize}
\\
		\hline
	\end{tabular}
	\egroup
\end{center}

\subsubsection{Comptes normaux}
\begin{center}
	\bgroup
	\def\arraystretch{1.5}
	\begin{tabular}{|p{2.5cm}|p{2cm}|p{8cm}|p{8cm}|p{5cm}|}
		\hline
		\rowcolor{gris}Référence & Exigence & Test & Résultat attendu & Résultat/Commentaires\\
		\hline
		ConS.1 & 3.1.2.6.6.1 & Ouvrir le dossier contenant les comptes normaux & Vérifier que le fichier à bien pour extension .acc &\vspace*{-0.2cm} \begin{itemize}[label=$\ast$] \item \textcolor{true}{Validé.}\end{itemize}
 \\
		\cline{1-2}\ccline{3-3}{gris}\cline{4-5}
		ConS.2 & 3.1.2.6.6.2 & Ouvrir le fichier .acc & Les lignes sont de la forme "pseudo@motDePasse:/fichier\#/autrefichier" &\vspace*{-0.2cm} \begin{itemize}[label=$\ast$] \item \textcolor{true}{Validé.}\end{itemize} \\
		\hline
	\end{tabular}
	\egroup
\end{center}

\subsubsection{Comptes temporaires}
\begin{center}
	\bgroup
	\def\arraystretch{1.5}
	\begin{tabular}{|p{2.5cm}|p{2cm}|p{8cm}|p{8cm}|p{5cm}|}
		\hline
		\rowcolor{gris}Référence & Exigence & Test & Résultat attendu & Résultat/Commentaires\\
		\hline
		CotS.1 & 3.1.2.6.7.1 & Ouvrir le dossier contenant les comptes temporaires. & L'extension du fichier est .acctmp &\vspace*{-0.2cm} \begin{itemize}[label=$\ast$] \item \textcolor{false}{Non impémenté}\end{itemize}\\
		\cline{1-2}\ccline{3-3}{gris}\cline{4-5}
		CotS.2 & 3.1.2.6.7.2 & Ouvrir le fichier .acctmp & Le fichier contient un timestamp. &\vspace*{-0.2cm} \begin{itemize}[label=$\ast$] \item \textcolor{false}{Non impémenté}\end{itemize}
\\
		\cline{1-2}\ccline{3-3}{gris}\cline{4-5}
		CotS.3 & 3.1.2.6.7.3 & Accéder à un compte dont le temps est écoulé & Le compte doit être supprimé du fichier .acctmp &\vspace*{-0.2cm} \begin{itemize}[label=$\ast$] \item \textcolor{false}{Non impémenté}\end{itemize}
\\
		\cline{1-2}\cline{4-5}
		CotS.3.1 & 3.1.2.6.7.4 & & Un message indique que le compte temporaire vient d'être supprimé. &\vspace*{-0.2cm} \begin{itemize}[label=$\ast$] \item \textcolor{false}{Non impémenté}\end{itemize}
\\
		\cline{1-2}\ccline{3-3}{gris}\cline{4-5}
		CotS.4 & 3.1.2.6.7.5 & Ouvrir le fichier .acctmp & Les lignes sont de la forme "DateFinEnSeconde|pseudo@motDePasse:/fichier\#/autrefichier" &\vspace*{-0.2cm} \begin{itemize}[label=$\ast$] \item \textcolor{false}{Non impémenté}\end{itemize}
\\
		\hline
	\end{tabular}
	\egroup
\end{center}

\subsubsection{Comptes bannis}
\begin{center}
	\bgroup
	\def\arraystretch{1.5}
	\begin{tabular}{|p{2.5cm}|p{2cm}|p{8cm}|p{8cm}|p{5cm}|}
		\hline
		\rowcolor{gris}Référence & Exigence & Test & Résultat attendu & Résultat/Commentaires\\
		\hline
		CobS.1 & 3.1.2.6.8.1 & Ouvrir le dossier contenant les comptes bannis & Le fichier contenant les comptes bannis doit avoir pour extension .accban &\vspace*{-0.2cm} \begin{itemize}[label=$\ast$] \item \textcolor{false}{Non impémenté}\end{itemize}\\
		\cline{1-2}\ccline{3-3}{gris}\cline{4-5}
		CobS.2 & 3.1.2.6.8.2 & Ouvrir le fichier .accban & Chaque ligne contient le motif de bannissement. &\vspace*{-0.2cm} \begin{itemize}[label=$\ast$] \item \textcolor{false}{Non impémenté}\end{itemize}
\\
		\cline{1-2}\cline{4-5}
		CobS.3 & 3.1.2.6.8.3 & & Un timestamp est présent dans le fichier. &\vspace*{-0.2cm} \begin{itemize}[label=$\ast$] \item \textcolor{false}{Non impémenté}\end{itemize}
\\
		\cline{1-2}\cline{4-5}
		CobS.4 & 3.1.2.6.8.4 & & Chaque ligne est de la forme "TimestampDeFin|pseudo@motDePasse:/fichier\#/autrefichier" &\vspace*{-0.2cm} \begin{itemize}[label=$\ast$] \item \textcolor{false}{Non impémenté}\end{itemize}
\\
		\cline{1-2}\cline{4-5}
		CobS.5 & 3.1.2.6.8.5 & & Quand le timestamp actuel est égale au timestamp présent dans le fichier, le compte repasse dans le fichier .acc &\vspace*{-0.2cm} \begin{itemize}[label=$\ast$] \item \textcolor{false}{Non impémenté}\end{itemize}
\\
		\hline
	\end{tabular}
	\egroup
\end{center}

\subsection{Limitations}
\begin{center}
	\bgroup
	\def\arraystretch{1.5}
	\begin{tabular}{|p{2.5cm}|p{2cm}|p{8cm}|p{8cm}|p{5cm}|}
		\hline
		\rowcolor{gris}Référence & Exigence & Test & Résultat attendu & Résultat/Commentaires\\
		\hline
		LimiS.1 & 3.1.2.7.1 & Ouvrir le programme & Vérifier la présence du menu "Configuration" dans la barre des menus. &\vspace*{-0.2cm} \begin{itemize}[label=$\ast$] \item \textcolor{false}{Non impémenté}\end{itemize}\\
		\cline{1-2}\ccline{3-3}{gris}\cline{4-5}
		LimiS.2 & 3.1.2.7.2 & Cliquer sur configuration, puis sur "limiter la vitesse" et mettre en place une limite & Le programme ne télécharge pas plus vite que la vitesse définie. &\vspace*{-0.2cm} \begin{itemize}[label=$\ast$] \item \textcolor{false}{Non impémenté}\end{itemize}
\\
		\cline{1-2}\ccline{3-3}{gris}\cline{4-5}
		LimiS.3 & 3.1.2.7.3 & Ouvrir le menu déroulant "Configuration" & La liste déroulante contient un bouton "Limiter vitesse de téléchargement" &\vspace*{-0.2cm} \begin{itemize}[label=$\ast$] \item \textcolor{false}{Non impémenté}\end{itemize}
\\
		\cline{1-2}\ccline{3-3}{gris}\cline{4-5}
		LimiS.4 & 3.1.2.7.4 & Cliquer sur configuration, puis sur "limiter la vitesse" & Un curseur précis apparait pour régler la vitesse de téléchargement. &\vspace*{-0.2cm} \begin{itemize}[label=$\ast$] \item \textcolor{false}{Non impémenté}\end{itemize}
\\
		\cline{1-2}\cline{4-5}
		LimiS.5 & 3.1.2.7.5 & & Un curseur apparait pour régler la vitesse de téléchargement. &\vspace*{-0.2cm} \begin{itemize}[label=$\ast$] \item \textcolor{false}{Non impémenté}\end{itemize}
\\
		\cline{1-2}\ccline{3-3}{gris}\cline{4-5}
		LimiS.6 & 3.1.2.7.6 & & La vitesse de téléchargement de l'utilisateur est limitée. &\vspace*{-0.2cm} \begin{itemize}[label=$\ast$] \item \textcolor{false}{Non impémenté}\end{itemize}
\\
		\cline{1-2}\ccline{3-3}{gris}\cline{4-5}
		LimiS.7 & 3.1.2.7.7 & Ouvrir le menu déroulant Configuration & Un bouton "Limiter vitesse de téléversement" est présent &\vspace*{-0.2cm} \begin{itemize}[label=$\ast$] \item \textcolor{false}{Non impémenté}\end{itemize}
\\
		\cline{1-2}\ccline{3-3}{gris}\cline{4-5}
		LimiS.8 & 3.1.2.7.8 & Cliquer sur "Limiter vitesse de téléversement" & Un curseur apparait. &\vspace*{-0.2cm} \begin{itemize}[label=$\ast$] \item \textcolor{false}{Non impémenté}\end{itemize}
\\
		\cline{1-2}\cline{4-5}
		LimiS.9 & 3.1.2.7.9 & & Le curseur permet de donner une valeur limite. &\vspace*{-0.2cm} \begin{itemize}[label=$\ast$] \item \textcolor{false}{Non impémenté}\end{itemize}
\\
		\cline{1-2}\ccline{3-3}{gris}\cline{4-5}
		LimiS.10 & 3.1.2.7.11 & Ouvrir le menu déroulant "Configuration" & Un bouton "Plages horaires téléchargement" est présent. &\vspace*{-0.2cm} \begin{itemize}[label=$\ast$] \item \textcolor{false}{Non impémenté}\end{itemize}
\\
		\cline{1-2}\ccline{3-3}{gris}\cline{4-5}
		LimiS.11 & 3.1.2.7.12 & Cliquer sur le bouton "Plages horaires téléchargement" & Une liste d’horaires à cocher apparait. &\vspace*{-0.2cm} \begin{itemize}[label=$\ast$] \item \textcolor{false}{Non impémenté}\end{itemize}
\\
		\cline{1-2}\ccline{3-3}{gris}\cline{4-5}
		LimiS.12 & 3.1.2.7.10 & Bloquer les téléchargement sur une plage horaire et tester de télécharger un fichier dans cette plage horaire & Les téléchargement sont bloqués sur la plage horaire définie. &\vspace*{-0.2cm} \begin{itemize}[label=$\ast$] \item \textcolor{false}{Non impémenté}\end{itemize}
\\
		\cline{1-2}\ccline{3-3}{gris}\cline{4-5}
		LimiS.13 & 3.1.2.7.14 & Ouvrir le menu "Configuration" & Un bouton "Plage horaire téléversement" est présent dans la liste déroulante. &\vspace*{-0.2cm} \begin{itemize}[label=$\ast$] \item \textcolor{false}{Non impémenté}\end{itemize}
\\
		\cline{1-2}\ccline{3-3}{gris}\cline{4-5}
		LimiS.14 & 3.1.2.7.15 & Cliquer sur le bouton "Plage horaire téléversement" & Une liste d’horaires à cocher apparait. &\vspace*{-0.2cm} \begin{itemize}[label=$\ast$] \item \textcolor{false}{Non impémenté}\end{itemize}
\\
		\cline{1-2}\ccline{3-3}{gris}\cline{4-5}
		LimiS.15 & 3.1.2.7.13 & Bloquer les téléchargement sur une plage horaire et tester de télécharger un fichier dans cette plage horaire & Les téléchargement sont bloqués sur la plage horaire définie. &\vspace*{-0.2cm} \begin{itemize}[label=$\ast$] \item \textcolor{false}{Non impémenté}\end{itemize}
\\
		\cline{1-2}\ccline{3-3}{gris}\cline{4-5}
		LimiS.16 & 3.1.2.7.16 & Essayer d'uploader un fichier dont l'extension est interdite ( exemple .exe) & Le programme doit refuser l'upload. &\vspace*{-0.2cm} \begin{itemize}[label=$\ast$] \item \textcolor{false}{Non impémenté}\end{itemize}
\\
		\cline{1-2}\ccline{3-3}{gris}\cline{4-5}
		LimiS.17 & 3.1.2.7.17 & Cliquer sur le menu "Configuration" & Un bouton "Extensions interdites" est présent dans la liste déroulante. &\vspace*{-0.2cm} \begin{itemize}[label=$\ast$] \item \textcolor{false}{Non impémenté}\end{itemize}
\\
		\hline
	\end{tabular}
	\egroup
\end{center}
\newpage
\begin{center}
	\bgroup
	\def\arraystretch{1.5}
	\begin{tabular}{|p{2.5cm}|p{2cm}|p{8cm}|p{8cm}|p{5cm}|}
		\hline
		\rowcolor{gris}Référence & Exigence & Test & Résultat attendu & Résultat/Commentaires\\
		\hline
		LimiS.18 & 3.1.2.7.18 & Cliquer sur le bouton "Extensions interdites" & Un champ de saisie apparait. &\vspace*{-0.2cm} \begin{itemize}[label=$\ast$] \item \textcolor{false}{Non impémenté}\end{itemize}
\\
		\cline{1-2}\cline{4-5}
		LimiS.19 & 3.1.2.7.19 & & Le champs est pré-rempli avec les extensions déjà interdite. &\vspace*{-0.2cm} \begin{itemize}[label=$\ast$] \item \textcolor{false}{Non impémenté}\end{itemize}
\\
		\cline{1-2}\cline{4-5}
		LimiS.20 & 3.1.2.7.20 & & Un bouton valide est présent en dessous du champs présent. &\vspace*{-0.2cm} \begin{itemize}[label=$\ast$] \item \textcolor{false}{Non impémenté}\end{itemize}
\\
		\hline
	\end{tabular}
	\egroup
\end{center}

\subsection{Utilisation des dossiers}
\begin{center}
	\bgroup
	\def\arraystretch{1.5}
	\begin{tabular}{|p{2.5cm}|p{2cm}|p{8cm}|p{8cm}|p{5cm}|}
		\hline
		\rowcolor{gris}Référence & Exigence & Test & Résultat attendu & Résultat/Commentaires\\
		\hline
		UtdoS.1 & 3.1.2.8.1 & Ouvrir le logiciel. & L'arborescence des dossiers sur le serveur est affichée sur la droite de la fenêtre. &\vspace*{-0.2cm} \begin{itemize}[label=$\ast$] \item \textcolor{false}{Non impémenté}\end{itemize}
\\
		\cline{1-2}\ccline{3-3}{gris}\cline{4-5}
		UtdoS.1.1 & 3.1.2.8.2 3.1.2.8.3 & Cliquer sur les dossiers dans l'arborescence. & Les sous-dossiers et fichiers des dossiers sur lesquels on a cliqué sont affichés dans l'arborescence. &\vspace*{-0.2cm} \begin{itemize}[label=$\ast$] \item \textcolor{false}{Non impémenté}\end{itemize}
\\
		\cline{1-2}\ccline{3-3}{gris}\cline{4-5}
		UtdoS.1.2 & 3.1.2.8.5 3.1.2.8.6 3.1.2.8.10 3.1.2.8.13 & Faire un clic droit sur un dossier. & Un menu de gestion de dossier apparait, les boutons "Nouveau sous-dossier", "Supprimer le dossier", et "Nouveau fichier vide" sont disponible à l'intérieur. &\vspace*{-0.2cm} \begin{itemize}[label=$\ast$] \item \textcolor{false}{Non impémenté}\end{itemize}
\\
		\cline{1-2}\ccline{3-3}{gris}\cline{4-5}
		UtdoS.1.2.1 & 3.1.2.8.7 3.1.2.8.8 & Cliquer sur "Nouveau sous-dossier". & Une boite de dialogue apparait avec un champ de saisie "nom". &\vspace*{-0.2cm} \begin{itemize}[label=$\ast$] \item \textcolor{false}{Non impémenté}\end{itemize}
\\
		\cline{1-2}\ccline{3-3}{gris}\cline{4-5}
		UtdoS.1.2.1.1 & 3.1.2.8.4 & Entrer un nom et valider la boite de dialogue. & Un nouveau sous-dossier avec le nom donné dans la boite de dialogue apparait dans l'arborescence. &\vspace*{-0.2cm} \begin{itemize}[label=$\ast$] \item \textcolor{false}{Non impémenté}\end{itemize}
\\
		\cline{1-2}\ccline{3-3}{gris}\cline{4-5}
		UtdoS.1.2.2 & 3.1.2.8.11 3.1.2.8.11.1 & Cliquer sur "Supprimer le dossier". & Une boite de dialogue proposant un bouton "Valider" et un "Annuler" apparait. &\vspace*{-0.2cm} \begin{itemize}[label=$\ast$] \item \textcolor{false}{Non impémenté}\end{itemize}
\\
		\cline{1-2}\ccline{3-3}{gris}\cline{4-5}
		UtdoS.1.2.2.1 & 3.1.2.8.9 3.1.2.8.11.2 & Cliquer sur "Valider". & La boite de dialogue disparait et le dossier est supprimé. &\vspace*{-0.2cm} \begin{itemize}[label=$\ast$] \item \textcolor{false}{Non impémenté}\end{itemize}
\\
		\cline{1-2}\ccline{3-3}{gris}\cline{4-5}
		UtdoS.1.2.2.2 & 3.1.2.8.11.3 & Cliquer sur "Annuler". & La boite de dialogue disparait sans rien déclencher. &\vspace*{-0.2cm} \begin{itemize}[label=$\ast$] \item \textcolor{false}{Non impémenté}\end{itemize}
\\
		\cline{1-2}\ccline{3-3}{gris}\cline{4-5}
		UtdoS.1.2.3 & 3.1.2.8.12 3.1.2.8.14 & Cliquer sur "Nouveau fichier vide". & Un nouveau fichier "nouveau\_fichier" apparait dans le dossier. &\vspace*{-0.2cm} \begin{itemize}[label=$\ast$] \item \textcolor{false}{Non impémenté}\end{itemize}
\\
		\cline{1-2}\ccline{3-3}{gris}\cline{4-5}
		UtdoS.1.3 & 3.1.2.8.16 3.1.2.8.17 & Faire un clic droit sur un des fichiers dans l'arborescence. & Un menu de gestion de fichier apparait avec un bouton "Supprimer le fichier" à l'intérieur. &\vspace*{-0.2cm} \begin{itemize}[label=$\ast$] \item \textcolor{false}{Non impémenté}\end{itemize}
\\
		\cline{1-2}\ccline{3-3}{gris}\cline{4-5}
		UtdoS.1.3.1 & 3.1.2.8.18 3.1.2.8.18.1 & Cliquer sur "Supprimer le fichier". & Une boite de dialogue proposant un bouton "Valider" et un "Annuler" apparait. &\vspace*{-0.2cm} \begin{itemize}[label=$\ast$] \item \textcolor{false}{Non impémenté}\end{itemize}
\\
		\cline{1-2}\ccline{3-3}{gris}\cline{4-5}
		UtdoS.1.3.1.1 & 3.1.2.8.15 3.1.2.8.18.2 & Cliquer sur "Valider". & La boite de dialogue disparait et le dossier est supprimé. &\vspace*{-0.2cm} \begin{itemize}[label=$\ast$] \item \textcolor{false}{Non impémenté}\end{itemize}
\\
		\cline{1-2}\ccline{3-3}{gris}\cline{4-5}
		UtdoS.1.3.1.2 & 3.1.2.8.18.3 & Cliquer sur "Annuler". & La boite de dialogue disparait sans rien déclencher. &\vspace*{-0.2cm} \begin{itemize}[label=$\ast$] \item \textcolor{false}{Non impémenté}\end{itemize}
\\
		\hline
	\end{tabular}
	\egroup
\end{center}

\subsection{Communication réseau}

\subsubsection{Procédure d’initialisation de la communication}
\begin{center}
	\bgroup
	\def\arraystretch{1.5}
	\begin{tabular}{|p{2.5cm}|p{2cm}|p{8cm}|p{8cm}|p{5cm}|}
		\hline
		\rowcolor{gris}Référence & Exigence & Test & Résultat attendu & Résultat/Commentaires\\
		\hline
		PicoS.1 & 3.1.2.9.1.1 & Après le démarrage du serveur, lancer la commande "netstat -a -p" sur le terminal du client & Le terminal doit indiquer que le port 21 est à l'état "LISTEN" et le nom ou le PID doit correspondre à celui du logiciel serveur &\vspace*{-0.2cm} \begin{itemize}[label=$\ast$] \item \textcolor{true}{Validé.}\end{itemize} \\
		\hline
		PicoS.2 & 3.1.2.9.1.1 & Après le lancement d'un transfert de fichier, lancer la commande "netstat -a -p" sur le terminal du serveur & Le terminal doit afficher une connexion avec une adresse IP correspondant au client, dont le port de connexion de l'adresse \textbf{locale} doit être différent de 21 (soit supérieur à 1024, soit égale à 20), l'état doit être à "ESTABLISHED" et le nom ou le PID doit être celui du logiciel serveur.&\vspace*{-0.2cm} \begin{itemize}[label=$\ast$] \item \textcolor{true}{Validé.}\end{itemize}\\
		\cline{1-2}
		PicoS.3 & 3.1.2.9.1.3 & & &\\
		\hline
	\end{tabular}
	\egroup
\end{center}

\subsubsection{Download et Upload}
\begin{center}
	\bgroup
	\def\arraystretch{1.5}
	\begin{tabular}{|p{2.5cm}|p{2cm}|p{8cm}|p{8cm}|p{5cm}|}
		\hline
		\rowcolor{gris}Référence & Exigence & Test & Résultat attendu & Résultat/Commentaires\\
		\hline
		DwupS.1 & 3.1.2.10.1 & Télécharger puis Téléverser un fichier en étant connecté & Le transférer se fait dans les 2 cas. &\vspace*{-0.2cm} \begin{itemize}[label=$\ast$] \item \textcolor{false}{Non impémenté}\end{itemize}
\\
		\hline
	\end{tabular}
	\egroup
\end{center}

\subsubsection{Module de statistique}
\begin{center}
	\bgroup
	\def\arraystretch{1.5}
	\begin{tabular}{|p{2.5cm}|p{2cm}|p{8cm}|p{8cm}|p{5cm}|}
		\hline
		\rowcolor{gris}Référence & Exigence & Test & Résultat attendu & Résultat/Commentaires\\
		\hline
		MostS.1 & 3.1.2.11.1.1 & Ouvrir le logiciel et cliquer sur le bouton statistique. & Le programme affiche la fréquence de connexion des utilisateurs. &\vspace*{-0.2cm} \begin{itemize}[label=$\ast$] \item \textcolor{false}{Non impémenté}\end{itemize}\\
		\cline{1-2}\cline{4-5}
		MostS.2 & 3.1.2.11.1.2 & & Le programme affiche le nombre de téléchargement de chaque fichier. &\vspace*{-0.2cm} \begin{itemize}[label=$\ast$] \item \textcolor{false}{Non impémenté}\end{itemize}
\\
		\cline{1-2}\cline{4-5}
		MostS.3 & 3.1.2.11.1.3 & & Le programme affiche le nombre de téléchargement par utilisateur. &\vspace*{-0.2cm} \begin{itemize}[label=$\ast$] \item \textcolor{false}{Non impémenté}\end{itemize}
\\
		\cline{1-2}\cline{4-5}
		MostS.4 & 3.1.2.11.1.4 & & Le programme affiche le nombre de téléversement effectué par utilisateur. &\vspace*{-0.2cm} \begin{itemize}[label=$\ast$] \item \textcolor{false}{Non impémenté}\end{itemize}
\\
		\cline{1-2}\cline{4-5}
		MostS.5 & 3.1.2.11.1.5 & & Le programme affiche le nombre de téléchargement par tranche de 2 heures. &\vspace*{-0.2cm} \begin{itemize}[label=$\ast$] \item \textcolor{false}{Non impémenté}\end{itemize}
\\
		\cline{1-2}\cline{4-5}
		MostS.6 & 3.1.2.11.1.6 & & Un bouton permet de paramétrer le nombre d'heure entre chaque tranches horaires. &\vspace*{-0.2cm} \begin{itemize}[label=$\ast$] \item \textcolor{false}{Non impémenté}\end{itemize}
\\
		\cline{1-2}\cline{4-5}
		MostS.7 & 3.1.2.11.1.7 & & La programme affiche en temps réel les statistiques suivante : le fichier le plus téléchargé, l'utilisateur qui upload le plus de fichier, l'utilisateur qui télécharge le plus de fichier, un histogramme traçant le nombre de téléchargement de chaque fichier, un histogramme du nombre de fichier téléchargé et téléversé par utilisateur, un histogramme du nombre moyen de téléchargement et téléversement par tranche de 2 heures. &\vspace*{-0.2cm} \begin{itemize}[label=$\ast$] \item \textcolor{false}{Non impémenté}\end{itemize}
\\
		\cline{1-2}\cline{4-5}
		MostS.7.1 & 3.1.2.11.1.7 & & Toutes les 2 semaines un rapport contient les informations listé ci-dessus. &\vspace*{-0.2cm} \begin{itemize}[label=$\ast$] \item \textcolor{false}{Non impémenté}\end{itemize}
\\
		\cline{1-2}\ccline{3-3}{gris}\cline{4-5}
		MostS.8 & 3.1.2.11.1.8 & Ouvrir le logiciel & Un bouton "Temps réel" est présent et génère un fichier. &\vspace*{-0.2cm} \begin{itemize}[label=$\ast$] \item \textcolor{false}{Non impémenté}\end{itemize}
\\
		\cline{1-2}\ccline{3-3}{gris}\cline{4-5}
		MostS.9 & 3.1.2.11.1.9 & & Le fichier généré aura comme format .html &\vspace*{-0.2cm} \begin{itemize}[label=$\ast$] \item \textcolor{false}{Non impémenté}\end{itemize}
\\
		\hline
	\end{tabular}
	\egroup
\end{center}

\subsubsection{Fichier d’historique des événements}
\begin{center}
	\bgroup
	\def\arraystretch{1.5}
	\begin{tabular}{|p{2.5cm}|p{2cm}|p{8cm}|p{8cm}|p{5cm}|}
		\hline
		\rowcolor{gris}Référence & Exigence & Test & Résultat attendu & Résultat/Commentaires\\
		\hline
		FhevS.1 & 3.1.2.12.1.1 & Lancer le serveur. & Un fichier d'historique d'événement est généré. &\vspace*{-0.2cm} \begin{itemize}[label=$\ast$] \item \textcolor{false}{Non impémenté}\end{itemize}\\
		\cline{1-2}\cline{4-5}
		FhevS.2 & 3.1.2.12.1.2 & & Le fichier contient tout les événement passé entre 0h00 et 23h 59m 59s&\vspace*{-0.2cm} \begin{itemize}[label=$\ast$] \item \textcolor{false}{Non impémenté}\end{itemize}
\\
		\cline{1-2}\cline{4-5}
		FhevS.3 & 3.1.2.12.1.3 & & Le fichier se trouve dans un répertoire log présent dans le répertoire du logiciel. &\vspace*{-0.2cm} \begin{itemize}[label=$\ast$] \item \textcolor{false}{Non impémenté}\end{itemize}
\\
		\cline{1-2}\ccline{3-3}{gris}\cline{4-5}
		FhevS.4 & 3.1.2.12.1.4 & Ouvrir le logiciel & Un bouton log est présent dans la barre d'outil. &\vspace*{-0.2cm} \begin{itemize}[label=$\ast$] \item \textcolor{false}{Non impémenté}\end{itemize}
\\
		\hline
	\end{tabular}
	\egroup
\end{center}

\subsection{Cryptage}

\begin{center}
	\bgroup
	\def\arraystretch{1.5}
	\begin{tabular}{|p{2.5cm}|p{2cm}|p{8cm}|p{8cm}|p{5cm}|}
		\hline
		\rowcolor{gris}Référence & Exigence & Test & Résultat attendu & Résultat/Commentaires\\
		\hline
		CrypS.1 & 3.1.3.1.1 & Ouvrir le logiciel & Vérifier la présence des boutons de cryptage dans le menu configuration. &\vspace*{-0.2cm} \begin{itemize}[label=$\ast$] \item \textcolor{false}{Non impémenté}\end{itemize}\\
		\cline{1-2}\ccline{3-3}{gris}\cline{4-5}
		CrypS.2 & 3.1.3.1.2 & Ourir le menu de gestion des fichiers. & Un bouton "Cryptage du fichier" est présent. &\vspace*{-0.2cm} \begin{itemize}[label=$\ast$] \item \textcolor{false}{Non impémenté}\end{itemize}
\\
		\cline{1-2}\ccline{3-3}{gris}\cline{4-5}
		CrypS.3 & 3.1.3.1.3 & Définir une clé de cryptage pour un fichier. & Le programme indique dans une console de log que le clé de cryptage est sauvegardée. &\vspace*{-0.2cm} \begin{itemize}[label=$\ast$] \item \textcolor{false}{Non impémenté}\end{itemize}
\\
		\cline{1-2}\ccline{3-3}{gris}\cline{4-5}
		CrypS.4 & 3.1.3.1.4 & Cliquer sur "Cryptage du fichier" et définir une clé & Le fichier est accessible uniquement avec la clé. &\vspace*{-0.2cm} \begin{itemize}[label=$\ast$] \item \textcolor{false}{Non impémenté}\end{itemize}
\\
		\cline{1-2}\ccline{3-3}{gris}\cline{4-5}
		CrypS.5 & 3.1.3.1.5 & Cliquer sur un fichier Crypté et entré la clé. & Le fichier est lisible et décrypté. &\vspace*{-0.2cm} \begin{itemize}[label=$\ast$] \item \textcolor{false}{Non impémenté}\end{itemize}
\\
		\cline{1-2}\ccline{3-3}{gris}\cline{4-5}
		CrypS.6 & 3.1.3.1.6 & Cliquer sur "décrypté le fichier." & Le programme décrypte le fichier avec la clé donnée. &\vspace*{-0.2cm} \begin{itemize}[label=$\ast$] \item \textcolor{false}{Non impémenté}\end{itemize}
\\
		\cline{1-2}\ccline{3-3}{gris}\cline{4-5}
		CrypS.7 & 3.1.3.1.7 & Téléchargé un fichier crypté. & Un champs de saisie apparait. &\vspace*{-0.2cm} \begin{itemize}[label=$\ast$] \item \textcolor{false}{Non impémenté}\end{itemize}
\\
		\cline{1-2}\cline{4-5}
		CrypS.8 & 3.1.3.1.8 & & Le champs peut être rempli à l'aide d'une clé. &\vspace*{-0.2cm} \begin{itemize}[label=$\ast$] \item \textcolor{false}{Non impémenté}\end{itemize}
\\
		\cline{1-2}\ccline{3-3}{gris}\cline{4-5}
		CrypS.9 & 3.1.3.1.9 & Entrer la clé de cryptage et valider. & Le programme indique que la clé à été enregistré dans une console de log. &\vspace*{-0.2cm} \begin{itemize}[label=$\ast$] \item \textcolor{false}{Non impémenté}\end{itemize}
\\
		\cline{1-2}\ccline{3-3}{gris}\cline{4-5}
		CrypS.10 & 3.1.3.1.10 & Crypter un fichier sur le serveur et le télécharger sur un client. & Le client entre une clé, si celle ci est en accord avec celle sur le serveur, le fichier est décrypté. &\vspace*{-0.2cm} \begin{itemize}[label=$\ast$] \item \textcolor{false}{Non impémenté}\end{itemize}
\\
		\cline{1-2}\ccline{3-3}{gris}\cline{4-5}
		CrypS.11 & 3.1.3.1.11 & Ouvrir le fichier contenant les clés. & Les clés ne sont pas écrit en clair. &\vspace*{-0.2cm} \begin{itemize}[label=$\ast$] \item \textcolor{false}{Non impémenté}\end{itemize}
\\
		\hline
	\end{tabular}
	\egroup
\end{center}

\end{document}